%%%%%%%%%%%%%%%%%%%%%%%%%%%%%%%%%%%%%%%%%%%%%%%%%%%%%%%%%%%%%%%%%%%%%%
% Lớp 11 | Dãy số. Cấp số cộng. Cấp số nhân | Cấp số nhân
% Dạng 1D2?3-2: Nhận diện cấp số nhân, công bội 𝑞
% Số câu: 116  (sinh tự động bởi tach_cau.py; từ nay thầy sửa trực tiếp file này)
%%%%%%%%%%%%%%%%%%%%%%%%%%%%%%%%%%%%%%%%%%%%%%%%%%%%%%%%%%%%%%%%%%%%%%

%%%%%%%%%%%%%%%% TRẮC NGHIỆM 4 PHƯƠNG ÁN %%%%%%%%%%%%%%%%

% nguồn: Ok_Toan11_TracNghiem_Lan2
\begin{ex}%[1D2N3-2]%[ma:46137007]
	Trong các dãy số sau, dãy  số nào là cấp số nhân?
	\choice 
		{$-3; 1; 5; 9; \ldots$} 
		{$\dfrac{1}{2}; \dfrac{2}{3}; \dfrac{3}{4};  \dfrac{4}{5}; \ldots$}
		{\True $16; 8; 4;  2; \ldots$}
		{$3; 6; 18; 108; \ldots$} 
	 \loigiai{
		\begin{itemize}
			\item Dãy số $-3; 1; 5; 9; \ldots$ không  là CSN vì $\dfrac{1}{-3}\neq \dfrac{5}{1}$.
			\item  Dãy số $\dfrac{1}{2}; \dfrac{2}{3}; \dfrac{3}{4}; \dfrac{4}{5}; \ldots$ không là CSN vì $\dfrac{2}{3}: \dfrac{1}{2}\neq \dfrac{3}{4}: \dfrac{2}{3}$.
			\item Dãy số $16; 8; 4; 2; \ldots$ là CSN vì $\forall n\in \mathbb{N}^*$  ta có $u_{n+1}=u_n\cdot q$.
			\item Dãy số $3; 6; 18; 108; \ldots$ không là CSN  vì $\dfrac{6}{3}\neq \dfrac{18}{6}$.
		\end{itemize}
	 }
\end{ex}

% nguồn: Ok_Toan11_TracNghiem_Lan2
\begin{ex}%[1D2N3-2]%[ma:451a0638]
	Cho cấp số nhân $\left(u_n\right)$, biết $u_1=1$,  $u_4=64$. Tính công bội $q$ của cấp số nhân đã cho
	\choice 
		{\True  $q=4$}
		{$q=-4$} 
		 {$q=21$}
		{$q=-2\sqrt{2}$} 
	\loigiai{Ta có $u_4=64 \Leftrightarrow u_1 q^3=64 \Leftrightarrow q^3=64  \Leftrightarrow q=4$.}
\end{ex}

% nguồn: Ok_Toan11_TracNghiem_Lan2
\begin{ex}%[1D2N3-2]%[ma:b22a3cef]
	Cho cấp số nhân $(u_n)$ biết $\heva{&u_4-u_2=54\\&u_5-u_3=108}$. Tìm số hạng đầu $u_1$ và công bội $q$ của  cấp số nhân trên.
	 \choice
		{\True $u_1=9$;  $q=2$}
		{$u_1=9$;  $q=-2$}
		{$u_1=-9$;  $q=-2$}
		{$u_1=-9$;  $q=2$}
	\loigiai{Ta có 
		\[ 
		\heva{&u_4-u_2=54\\&u_5-u_3=108} 
		 \Leftrightarrow \heva{&u_1q^3-u_1q=54\\&u_1q^4-u_1q^2=108}
		\Leftrightarrow \heva{&u_1q(q^2-1)=54\\&u_1q^2(q^2-1)=108} 
		\Leftrightarrow \heva{&u_1=9\\&q=2.} 
		 \]
	 }
\end{ex}

% nguồn: Ok_Toan11_TracNghiem_Lan2
\begin{ex}%[1D2N3-2]%[ma:d5a4af56]
	Tìm công bội $q$ của một cấp số nhân $(u_n)$ có $u_1=\dfrac{1}{2}$  và $u_6=16$.
	\choice 
		{$q=\dfrac{1}{2}$} 
		{$q=-2$} 
		 {\True $q=2$}
		{$q=-\dfrac{1}{2}$} 
	 \loigiai{
	Ta có $u_6=u_1q^5\Rightarrow  16=\dfrac{1}{2}\cdot q^5\Rightarrow q=2$.
	 }
\end{ex}

% nguồn: Ok_Toan11_TracNghiem_Lan2
\begin{ex}%[1D2N3-2]%[ma:4510ad76]
	Cho cấp số nhân $(u_n)$ với $u_1=-\dfrac{1}{2}$; $u_n=-32$. Tìm  $q$?
	\choice 
		 {$q=\pm \dfrac{1}{2}$}
		{\True $q=\pm 2$} 
		{$q=\pm 4$} 
		 {$q=\pm 1$}
	\loigiai{Áp dụng  công thức số hạng tổng quát cấp số nhân ta có\\
		$u_n = u_1 q^{n-1} \Leftrightarrow -32 = -\dfrac{1}{2} q^{n-1} \Leftrightarrow q^{n-1} = 64 \Leftrightarrow \left[\begin{array}{l} q = 2 \\ q = -2.
\end{array} \right.$}
\end{ex}

% nguồn: Ok_Toan11_TracNghiem_Lan2
\begin{ex}%[1D2N3-2]%[ma:bfab2814]
	Cho cấp số nhân $\left(u_n\right)$ với $u_1=3$ và $u_2=9$. Công bội của cấp số nhân  đã cho bằng
	 \choice
		 {$-6$}
		 {$\dfrac{1}{3}$}
		{\True $3$} 
		 {$6$}
	 \loigiai{
	 Ta có $u_2=u_1 \cdot q$ $\Rightarrow $ $q=\dfrac{u_2}{u_1}=3$.}
\end{ex}

% nguồn: Ok_Toan11_TracNghiem_Lan2
\begin{ex}%[1D2N3-2]%[ma:6ea7b7d5]
	Cho cấp số nhân $(u_n)$ có số  hạng đầu $u_1=-2$ và $u_6=486$. Công bội $q$ bằng
	 \choice
		{\True  $q=-3$}
		 {$q=5$}
		 {$q=\dfrac{3}{2}$}
		 {$q=-\dfrac{2}{3}$}
	\loigiai{Theo đề  ra ta có:$\heva{&u_1=-2\\&u_6=486}
		\Leftrightarrow  \heva{&u_1 = -2\\&486=u_1 q^5}
	\Rightarrow q^5 = -243  \Rightarrow q = -3$.}
\end{ex}

% nguồn: Ok_Toan11_TracNghiem_Lan2
\begin{ex}%[1D2N3-2]%[ma:5cefbef5]
	Cho cấp số nhân $\left(u_n\right)$ với $u_1=2$ và  công bội $q=3$. Tìm số hạng thứ $4$ của cấp số nhân.
	 \choice
		{$24$} 
		{\True  $54$}
		{$162$} 
		 {$48$}
	\loigiai{Có $u_4=u_1  q^3=2.3^3=54$.}
\end{ex}

% nguồn: Ok_Toan11_TracNghiem_Lan2
\begin{ex}%[1D2N3-2]%[ma:048cc934]
	Cho cấp số nhân $\left(u_n\right)$ với $u_1=3$ và $u_2=12$. Công bội của cấp số nhân đã cho bằng 
	 \choice
		 {$9$}
		{$-9$} 
		 {$\dfrac{1}{4}$}
		 {\True $4$}
	\loigiai{ 
	Công bội của cấp số nhân đã cho là $q=\dfrac{u_2}{u_1}=\dfrac{12}{3}=4$.} 
\end{ex}

% nguồn: Ok_Toan11_TracNghiem_Lan2
\begin{ex}%[1D2N3-2]%[ma:6590bdb7]
	Cho cấp số nhân $(u_n)$ có $u_4=9$, $u_5=81$. Công bội của cấp số  nhân đã cho là
	 \choice
		 {$3$}
		{$72$} 
		 {$18$}
		 {\True $9$}
	 \loigiai{
	Ta có công bội  $q=\dfrac{u_5}{u_4}=\dfrac{81}{9}=9$.
	 }
\end{ex}

% nguồn: Ok_Toan11_TracNghiem_Lan2
\begin{ex}%[1D2N3-2]%[ma:c56956eb]
	Dãy số  nào sau đây \textbf{không phải} là cấp số nhân?
	\choice 
		{\True  $1; 2; 3; 4; 5$}
		{$1;  2; 4; 8; 16$}
		{$1; -1; 1; -1;  1$}
		{$1; -2; 4;  -8; 16$}
	\loigiai{ 
		\begin{itemize}
			\item Dãy $1; 2; 4; 8; 16$ là cấp số nhân với công  bội $q=2$.
			\item Dãy $1; -1;  1; -1; 1$ là cấp số nhân với công bội $q=-1$.
			\item Dãy $1; -2; 4; -8; 16$ là cấp số nhân với  công bội $q=-2$.
			\item Dãy $1; 2; 3; 4; 5$ là cấp  số cộng với công sai $d=1$.
		\end{itemize}
	} 
\end{ex}

% nguồn: Ok_Toan11_TracNghiem_Lan2
\begin{ex}%[1D2N3-2]%[ma:815e2ae4]
	Dãy số nào sau đây \textbf{không phải} là cấp  số nhân?
	\choice 
		 {$1; -1; 1; -1$}
		{\True $1; -3;  9; 10$}
		{$1; 0; 0;  0$}
		{$32; 16;  8; 4$}
	\loigiai{ 
		Nếu  $(u_n)$ là cấp số nhân với công bội $q$ ta có
		$u_{n+1}=u_n\cdot q\Rightarrow  q=\dfrac{u_{n+1}}{u_n}$.
		\begin{itemize}
			\item $1; -1; 1;  -1$ là cấp số nhân với $q=-1$.
			\item $-1; 3; 9; 10$ không là cấp số  nhân.
			\item $1; 0; 0; 0$ là  cấp số nhân với $q=0$.
			\item $32; 16; 8; 4$ là cấp số  nhân với $q=\dfrac{1}{2}$.
		\end{itemize}
	 }
\end{ex}

% nguồn: Ok_Toan11_TracNghiem_Lan2
\begin{ex}%[1D2N3-2]%[ma:c4596227]
	Trong các dãy số sau, dãy số nào là một  cấp số nhân?
	\choice 
		{$1;  2; 3; 4; 5; 6; \ldots$}
		{$2; 4; 6; 8;  16; 32; \ldots$}
		{$-2;  -3; -4; -5; -6; -7; \ldots$}
		{\True $1; 2; 4; 8; 16;  32; \ldots$}
	 \loigiai{
		Dãy $1, 2; 4; 8; 16; 32; \ldots$  thỏa mãn: $u_{n+1}=2\cdot{u_n}\forall n\in \mathbb{N}^*$.\\
		Vậy dãy số $1; 2; 4; 8; 16; 32; \ldots$ là cấp số nhân với $u_1=1$ và công  bội $q=2$.
	 }
\end{ex}

% nguồn: Ok_Toan11_TracNghiem_Lan2
\begin{ex}%[1D2N3-2]%[ma:83d2ef99]
	Trong các dãy số  hữu hạn sau, dãy nào là cấp số nhân?
	 \choice
		 {$3; 9; 27; 54; 162$}
		{$81; 27; -9; -3; 1$} 
		 {\True $81; -27; 9; -3; 1$}
		{$3; 6;  9; 12; 15$}
	\loigiai{ 
		Xét các thương của  các số hạng với số hạng ngay trước nó, ta thấy:
		$\dfrac{-27}{81}=\dfrac{9}{-27}=\dfrac{-3}{9}=\dfrac{1}{-3}=-\dfrac{1}{3}$.  \\
		Vậy dãy $81; -27; 9; -3; 1$ là  cấp số nhân với công bội $q=-\dfrac{1}{3}$.
	} 
\end{ex}

% nguồn: Ok_Toan11_TracNghiem_Lan2
\begin{ex}%[1D2N3-2]%[ma:21b49699]
	Trong các dãy  số cho dưới đây, dãy số nào là cấp số nhân?
	\choice 
		 {$1; 2; 3; 4; 5$}
		{$1; 3;  6; 9; 12$}
		{$2;  4; 6; 8; 10$}
		{\True $2; 2; 2;  2; 2$}
	\loigiai{ 
		Ta thấy ở đáp án  $D$ có $u_1=u_2=u_3=u_4=u_5=2$ nên đây là cấp số nhân với công bội $q=1$.
	 }
\end{ex}

% nguồn: Ok_Toan11_TracNghiem_Lan2
\begin{ex}%[1D2N3-2]%[ma:d0ce2034]
	Tìm công bội $q$ của một  cấp số nhân $\left(u_n\right)$ có $u_1=\dfrac{1}{2}$ và $u_6=16$.
	 \choice
		 {$q=-\dfrac{1}{2}$}
		{\True $q=2$} 
		{$q=-2$} 
		 {$q=\dfrac{1}{2}$}
	\loigiai{Ta có  $u_6=u_1 q^5=16 \Rightarrow \dfrac{1}{2} q^5=16 \Leftrightarrow q^5=32 \Leftrightarrow q= 2$.}
\end{ex}

% nguồn: Ok_Toan11_TracNghiem_Lan2
\begin{ex}%[1D2N3-2]%[ma:902e9f76]
	Cho cấp số nhân $\left(u_n\right)$ với $u_1=2$ và $u_2=6$. Công bội của cấp số nhân đã cho bằng 
	\choice 
		{\True  $3$}
		 {$-4$}
		 {$4$}
		 {$\dfrac{1}{3}$}
	\loigiai{ 
	Công bội của cấp số nhân là  $q=\dfrac{u_2}{u_1}=\dfrac{6}{2}=3$.}
\end{ex}

% nguồn: Ok_Toan11_TracNghiem_Lan2
\begin{ex}%[1D2N3-2]%[ma:9e2279b6]
	Cấp số  nhân $\left(u_n\right)$ có $u_4=9$, $u_5=81$ có công bội là
	 \choice
		{$-3$} 
		 {$72$}
		{$18$} 
		{\True $9$} 
	\loigiai{Ta  có công bội $q=\dfrac{u_5}{u_4}=\dfrac{81}{9}=9$.}
\end{ex}

% nguồn: Ok_Toan11_TracNghiem_Lan2
\begin{ex}%[1D2N3-2]%[ma:1cf5307e]
	Biết ba số $x^2$; $8$; $x$ theo thứ tự lập thành cấp số nhân.  Giá trị của $x$ bằng
	 \choice
		 {$x=4$}
		{$x=5$} 
		{\True $x=2$} 
		 {$x=1$}
	\loigiai{Do ba số $x^2$; $8$; $x$ theo thứ tự lập thành cấp số nhân nên theo tính  chất cấp số nhân ta được
		\[x^2\cdot x  = 8
		\Leftrightarrow x^3 = 8 \Leftrightarrow x = 2.\] 
	} 
\end{ex}

% nguồn: Ok_Toan11_TracNghiem_Lan2
\begin{ex}%[1D2N3-2]%[ma:2cf24826]
	Cho cấp số nhân $(u_n)$ có $u_1=-\dfrac{1}{2}$, $u_7=-32$. Tìm công bội $q$. 
	\choice 
		 {$q=\pm\dfrac{1}{2}$}
		{\True $q=\pm 2$} 
		 {$q=\pm 4$}
		 {$q=\pm 1$}
	\loigiai{ 
	Ta  có $u_7=u_1q^6\Rightarrow -32=-\dfrac{1}{2}\cdot q^6\Rightarrow q=\pm 2$.
	 }
\end{ex}

% nguồn: Ok_Toan11_TracNghiem_Lan2
\begin{ex}%[1D2N3-2]%[ma:4dcc57e6]
	Dãy số nào sau đây \textbf{không phải} là cấp  số nhân?
	 \choice
		{$1; 2;  4; 8; \ldots$}
		 {$3; 3^2; 3^3; 3^4; \cdot \cdot$}
		{$4; 2; \text{; 1; }\dfrac{1}{2}; \ldots$} 
		{\True $\dfrac{1}{\pi}; \dfrac{1}{\pi^2}; \dfrac{1}{\pi^4};  \dfrac{1}{\pi^6}; \ldots$}
	 \loigiai{
		\begin{itemize}
			\item Dãy số $1; 2; 4; 8; \ldots$ là CSN với $u_1=1;  q=2$.
			\item Dãy số $3; 3^2; 3^3; 3^4;  \ldots$ là CSN với $u_1=3; q=3$.
			\item Dãy số $4; 2; 1; \dfrac{1}{2}; \ldots$ là CSN với $u_1=4; q=\dfrac{1}{2}$. 
			\item Dãy số $\dfrac{1}{\pi}; \dfrac{1}{\pi^2}; \dfrac{1}{\pi^4};  \dfrac{1}{\pi^6}; \ldots$ không phải CSN vì $\dfrac{u_2}{u_1}=\dfrac{1}{\pi}\neq q \dfrac{1}{\pi^2}=\dfrac{u_3}{u_2}$.
		\end{itemize}
	} 
\end{ex}

% nguồn: Ok_Toan11_TracNghiem_Lan2
\begin{ex}%[1D2N3-2]%[ma:90ef5cbd]
	Dãy số  nào sau đây \textbf{không phải} là cấp số nhân?
	\choice 
		{\True $1; -3;  9; -27; 54$}
		{$1; 2; 4; 8; 16$} 
		{$1; -1; 1; -1;  1$}
		{$1; -2; 4;  -8; 16$}
	\loigiai{ 
		\begin{itemize}
			\item Dãy $1; 2; 4; 8; 16$ là cấp số nhân với công bội $q=2$. 
			\item Dãy $1; -1; 1; -1; 1$  là cấp số nhân với công bội $q=-1$.
			\item Dãy $1; -2; 4; -8; 16$ là  cấp số nhân với công bội $q=-2$.
			\item Dãy $1; -3; 9;  -27; 54$ không phải là cấp số nhân vì $-3=1\cdot (-3); (-27)\cdot (-3)=81\neq 54$.
		\end{itemize}
	 }
\end{ex}

% nguồn: Ok_Toan11_TracNghiem_Lan2
\begin{ex}%[1D2N3-2]%[ma:a1bc7e46]
	Trong các dãy số sau dãy nào lập thành một cấp số nhân? 
	\choice 
		{$1; 3; 5;  7; 9$}
		{$1; 2;  4; 6; 8$}
		{$4;  \dfrac{1}{4}; 3; \dfrac{1}{3}; 2; \dfrac{1}{2}$}
		{\True $9; 3; 1;  \dfrac{1}{3}; \dfrac{1}{9}$}
	 \loigiai{
		Dãy các số $9; 3; 1; \dfrac{1}{3}; \dfrac{1}{9}$ là $5$ số hạng của $1$ cấp số nhân công bội  $q=\dfrac{1}{3}$.
	 }
\end{ex}

% nguồn: Ok_Toan11_TracNghiem_Lan2
\begin{ex}%[1D2N3-2]%[ma:2d5d71de]
	Cho cấp số nhân $\left(u_n\right)$ có $u_1=2$ và $u_4=54$.  Giá trị của công bội $q$ bằng
	 \choice
		{\True $3$} 
		 {$9$}
		{$27$} 
		 {$-3$}
	\loigiai{Ta có:$\dfrac{u_4}{u_1}=q^3 \Rightarrow  q^3=\dfrac{54}{2}=27 \Rightarrow q=\sqrt[3]{27}=3$.}
\end{ex}

% nguồn: Ok_Toan11_TracNghiem_Lan2
\begin{ex}%[1D2N3-2]%[ma:d292a585]
	Dãy số nào sau đây \textbf{không phải} là cấp số nhân? 
	\choice 
		{$4; -2;  1; \dfrac{-1}{2}$}
		{$1; 0; 0;  0$}
		{\True $1;  2; 4; 6$}
		 {$1; 1; 1; 1$}
	 \loigiai{
		Xét cấp số nhân $(u_n)$ với công bội $q$, ta có $u_n=u_{n-1}\cdot q$ với  $n\ge 2$.
		\begin{itemize}
			\item $4; -2; 1; \dfrac{-1}{2}$  là cấp số nhân với công bội $q=\dfrac{-1}{2}$.
			\item $1; 0; 0; 0$ là cấp số nhân  với công bội $q=0$.
			\item  $1; 2; 4; 6$ không là cấp số nhân.
			\item $1; 1; 1;  1$ là cấp số nhân với công bội $q=1$.
		\end{itemize}
	} 
\end{ex}

% nguồn: Ok_Toan11_TracNghiem_Lan2
\begin{ex}%[1D2N3-2]%[ma:7d5e1833]
	Cho cấp số nhân $(u_n)$ với $u_1=1$ và $u_2=2$.  Công bội của cấp số nhân đã cho là
	 \choice
		{$q=\dfrac{1}{2}$} 
		{\True  $q=2$}
		{$q=-2$} 
		{$q=-\dfrac{1}{2}$} 
	\loigiai{ 
	Ta có $u_2=u_1 \cdot q\Rightarrow q=\dfrac{u_2}{u_1}=\dfrac{2}{1}=2$.} 
\end{ex}

% nguồn: Ok_Toan11_TracNghiem_Lan2
\begin{ex}%[1D2N3-2]%[ma:de92c2c2]
	Cho cấp số nhân $(u_n)$ với $u_1=3$ và $u_2=15$. Công bội của cấp số nhân đã  cho bằng
	\choice 
		 {$-12$}
		 {$\dfrac{1}{5}$}
		{\True  $5$}
		{$12$} 
	\loigiai{ 
	Từ  định nghĩa cấp số nhân ta có $q=\dfrac{u_2}{u_1}=5$.}
\end{ex}

% nguồn: Ok_Toan11_TracNghiem_Lan2
\begin{ex}%[1D2N3-2]%[ma:0e797189]
	Trong các dãy số sau, dãy  số nào là cấp số nhân?
	\choice 
		{$3; 4;  5; 6; 7$}
		{$12; 14; 16; 18;  20$}
		{\True $3; 6; 12; 24;  48$}
		{$1;  1; 1; -1; 1$}
	 \loigiai{
		Vì $\dfrac{6}{3}=\dfrac{12}{6}=\dfrac{24}{12}=\dfrac{48}{24}(=2)$, nên dãy số $3; 6; 12;  24; 48$ là một cấp số nhân.
	 }
\end{ex}

% nguồn: Ok_Toan11_TracNghiem_Lan2
\begin{ex}%[1D2N3-2]%[ma:2b6258cd]
	Cho cấp số nhân $\left(u_n\right)$ có $u_1=-2$  và $u_5=-162$. Công bội $q$ bằng:
	\choice 
		{$q=-3$} 
		 {$q=3$}
		{\True $q=3;q=-3$} 
		 {$q=-2$}
	\loigiai{Ta có  $u_5=-162 \Leftrightarrow u_1 q^4=-162 \Leftrightarrow q^4=\dfrac{-162}{u_1}=\dfrac{-162}{-2}=81 \Leftrightarrow q=\pm 3$.}
\end{ex}

% nguồn: Ok_Toan11_TracNghiem_Lan2
\begin{ex}%[1D2N3-2]%[ma:68593e08]
	Cho cấp số nhân  $(u_n)$ có số hạng đầu $u_1=2$ và $u_6=486$. Công bội $q$ bằng
	 \choice
		{\True $q=3$} 
		 {$q=5$}
		{$q=\dfrac{3}{2}$} 
		{$q=\dfrac{2}{3}$} 
	\loigiai{ 
	 Ta có $u_6=u_1q^5\Rightarrow 486=2\cdot q^5\Rightarrow q=3$.
	 }
\end{ex}

% nguồn: Ok_Toan11_TracNghiem_Lan2
\begin{ex}%[1D2N3-2]%[ma:4df28690]
	Dãy nào sau đây là một cấp số  nhân?
	\choice 
		{$1; 2; 3;  4; \ldots$}
		{$1; 3;  5; 7; \ldots$}
		{\True $2;  4; 8; 16; \ldots$}
		{$2; 4; 6;  8; \ldots$}
	 \loigiai{
		Ta có $2; 4; 8; 16; \ldots$ là cấp  số nhân có số hạng đầu $u_1=2$ và công bội $q=2$.
	 }
\end{ex}

% nguồn: Ok_Toan11_TracNghiem_Lan2
\begin{ex}%[1D2N3-2]%[ma:cd56d3e9]
	Cho dãy số $(u_n)$ xác định  bởi $\heva{&u_1=-2\\&u_{n+1}=-\dfrac{1}{10}u_n}$. Chọn hệ thức đúng
	 \choice
		{\True  $(u_n)$ là cấp số nhân có công bội $q=-\dfrac{1}{10}$}
		{$u_n=(-2) \left(\dfrac{1}{10}\right)^{n-1}$} 
		{$u_n=\dfrac{u_{n-1}+u_{n+1}}{2}~(n  \geq 2) $}
		{$u_n=\sqrt{u_{n-1}u_{n+1}}~(n \geq 2)  $}
	\loigiai{Ta có $\dfrac{u_{n+1}}{u_n} = -\dfrac{1}{10}$  nên $(u_n)$ là cấp số nhân có công bội $q = -\dfrac{1}{10}$.}
\end{ex}

% nguồn: Ok_Toan11_TracNghiem_Lan2
\begin{ex}%[1D2N3-2]%[ma:89daafdd]
	Cho cấp số nhân $(u_n)$ có công bội dương và  $u_2=\dfrac{1}{4}$, $u_4=4$. Giá trị của $u_1$ là
	 \choice
		 {$u_1=-\dfrac{1}{6}$}
		{\True $u_1=\dfrac{1}{16}$} 
		{$u_1=-\dfrac{1}{16}$} 
		{$u_1=\dfrac{1}{2}$} 
	 \loigiai{Theo tính chất của cấp số nhân với $k \ge 2$ thì $u_k^2=u_{k-1}u_{k+1}$ ta suy ra\\
		\[u_3^2=u_2u_4=\dfrac{1}{4}.4=1 \Leftrightarrow \left[\begin{array}{l} u_3=1 \\ u_3=-1
	\end{array} \right.\]
	Vì $(u_n)$ là cấp số nhân có công bội dương nên  $u_3=1$.\\
Gọi $q$ là công bội ta được  $q=\dfrac{u_4}{u_3}=\dfrac{4}{1}=4$. Từ đó ta có $u_1=\dfrac{u_2}{q}=\dfrac{\dfrac{1}{4}}{4}=\dfrac{1}{16}$.}
\end{ex}

% nguồn: Ok_Toan11_TracNghiem_Lan2
\begin{ex}%[1D2N3-2]%[ma:90c916cf]
	Cho cấp số nhân  có $u_2=\dfrac{1}{4}$;$u_5=-16$. Tìm $q$ và $u_1$.
	 \choice
		{$q=-\dfrac{1}{2}$;  $u_1=-\dfrac{1}{2}$}
		{\True  $q=-4$; $u_1=-\dfrac{1}{16}$}
		{$q=\dfrac{1}{2}$;  $u_1=-\dfrac{1}{2}$}
		{$q=-4$; $u_1=\dfrac{1}{16}$} 
	\loigiai{Ta có $u_2 = u_1 q \Leftrightarrow \dfrac{1}{4} = u_1  q$; $u_5 = u_1 q^4 \Leftrightarrow -16 = u_1 q^4$.\\
	Suy ra $q^3 = -64  \Leftrightarrow q = -4$ từ đó $u_1 = -\dfrac{1}{16}$.}
\end{ex}

% nguồn: Ok_Toan11_TracNghiem_Lan2
\begin{ex}%[1D2N3-2]%[ma:5a93b914]
	Trong các dãy số sau, dãy số nào \textbf{không  phải} là một cấp số nhân?
	\choice 
		{$2; 4; 8; 16; \ldots$} 
		 {$1; -1; 1; -1; \ldots$}
		{\True $1^2; 2^2; 3^2; 4^2; \ldots$} 
		{$a; a^3;  a^5; a^7; \cdots (a\neq 0)$}
	\loigiai{ 
		Xét dãy số $1^2; 2^2; 3^2; 4^2; \ldots$ có $\dfrac{u_2}{u_1}=4\neq  q \dfrac{9}{4}=\dfrac{u_3}{u_2}$. \\
		Vậy dãy số $1^2; 2^2; 3^2; 4^2; \ldots$ không là CSN. 
		\begin{nx}
			Dãy $(u_n)$ với $u_n\neq 0$ là cấp số nhân khi $u_n=a\cdot{q^n}$, tức là các số hạng của nó đều được  biểu diễn dưới dạng lũy thừa của cùng một cơ số $q$, các số hạng liên tiếp thì số mũ của chúng cách đều nhau. Ví dụ
			\begin{itemize}
				\item $2; 4; 8; 16; \ldots$ là cấp số nhân và $u_n=2^n$. 
				\item $1; -1; 1;  -1; \ldots$ là cấp số nhân và $u_n=(-1)^n$.
				\item $a; a^3; a^5; a^7; \cdots (a\neq 0)$ là cấp số nhân và $u_n=a^{2n-1}=\dfrac{1}{a}\cdot (a^2)^n$. 
			\end{itemize}
		\end{nx}
	} 
\end{ex}

% nguồn: Ok_TracNghiemLop11_Lan1
\begin{ex}%[1D2N3-2]%[ma:f454dc05]
	Trong các dãy số cho bởi  công thức truy hồi sau, dãy số nào là cấp số nhân?
	 \choice
		 {$u_1=-1$, $u_{n+1}=u_n^2$}
		{\True $u_1=-1$, $u_{n+1}=2u_n$} 
		{$u_1=-1$,  $u_{n+1}=u_n+2$}
		{$u_1=-1$,  $u_{n+1}=u_n-2$}
	\loigiai{Theo  định nghĩa của cấp số nhân ta có dãy số cho bởi $u_1=-1$, $u_{n+1}=2 u_n$ là cấp số nhân.}
\end{ex}

% nguồn: Ok_TracNghiemLop11_Lan1
\begin{ex}%[1D2N3-2]%[ma:b3ac7c8e]
	Trong các  dãy số cho dưới đây, dãy số nào là cấp số nhân?
	 \choice
		{$1$; $2$; $3$; $4$; $5$} 
		{$0$;  $4$; $8$; $12$; $16$}
		{$1$; $3$; $5$; $7$; $9$} 
		{\True $ 5$;  $\dfrac{5}{2}$; $\dfrac{5}{4}$; $\dfrac{5}{8}$; $\dfrac{5}{16}$}
	\loigiai{ 
	Ta thấy ở dãy số $ 5$; $\dfrac{5}{2}$; $\dfrac{5}{4}$; $\dfrac{5}{8}$; $\dfrac{5}{16}$  có $u_1=5$, $u_2=\dfrac{5}{2}$, $u_3=\dfrac{5}{4}$, $u_4=\dfrac{5}{8}$, $u_5=\dfrac{5}{16}$ nên đây là cấp số nhân với công bội $ q=\dfrac{1}{2}$.}
\end{ex}

% nguồn: Ok_TracNghiemLop11_Lan1
\begin{ex}%[1D2N3-2]%[ma:4479efd0]
	Cho cấp số nhân $\left( u_n\right) $ có ${u_7=6}$ và $u_8=12$.  Tìm công bội $q$ cấp số nhân đã cho.
	\choice 
		 { $q=-6$}
		{ $q=6$} 
		 { $q=\dfrac{1}{2}$}
		{  \True $q=2$}
	 \loigiai{
		Ta  có $q=\dfrac{u_8}{u_7}=2$.
	} 
\end{ex}

% nguồn: Ok_TracNghiemLop11_Lan1
\begin{ex}%[1D2N3-2]%[ma:714e9066]
	Cho cấp số nhân $\left(u_n\right)$ có $u_1=1$, $u_2=4$. Công bội của cấp số nhân đã cho bằng 
	\choice 
		 {$21$}
		 {$-4$}
		{\True $4$} 
		{$2\sqrt{2}$} 
	 \loigiai{
		Công bội của cấp số  nhân là $q=\dfrac{u_2}{u_1}=4$.
	 }
\end{ex}

% nguồn: Ok_TracNghiemLop11_Lan1
\begin{ex}%[1D2N3-2]%[ma:b339e05c]
	Trong  các dãy số $\left(u_n\right)$ sau, dãy số nào là cấp số nhân?
	\choice  {$u_n=(-1)^n n$}{$u_n=n^2$}{$u_n=\dfrac{n}{3^n}$}{\True $u_n=2^n$}
	 \loigiai{
		\begin{itemize}
			\item $\dfrac{u_{n+1}}{u_n}=\dfrac{(-1)^{n+1} \cdot(n+1)}{(-1)^n \cdot  n}=-\dfrac{(n+1)}{n} \Rightarrow u_n=(-1)^n . n$ không phải cấp số nhân.
			\item $\dfrac{u_{n+1}}{u_n}=\dfrac{(n+1)^2}{n^2} \Rightarrow u_n=n^2$ không  phải cấp số nhân.
			\item $\dfrac{u_{n+1}}{u_n}=\dfrac{n+1}{3 n} \Rightarrow u_n=\dfrac{n}{3^n}$ không  phải cấp số nhân.
			\item $\dfrac{u_{n+1}}{u_n}=\dfrac{2^{n+1}}{2^n}=2 \Leftrightarrow u_{n+1}=2 u_n \Rightarrow u_n=2^n$ là  cấp số nhân.
		\end{itemize}
	} 
\end{ex}

% nguồn: Ok_TracNghiemLop11_Lan1
\begin{ex}%[1D2N3-2]%[ma:be64a514]
	Trong các dãy  số sau, dãy số nào \textbf{không} phải là một cấp số nhân?
	\choice 
		{$2;\, 4;\, 8;\,  16; \ldots$}
		{$1;\,-1;\,1;\,-1; \ldots$} 
		{\True $1^2;\, 2^2;\,  3^2;\, 4^2; \ldots$}
		{$a;\, a^3;\, a^5;\, a^7;  \ldots(a\ne 0)$}
	\loigiai{ 
		Các dãy  $\heva{&2;\, 4;\, 8;\, 16; \ldots\\
		&1;\,-1;\,1;\,-1;  \ldots\\
		&a;\, a^3;\, a^5;\, a^7;\, \ldots(a\ne 0)}$ đều là  các cấp số nhân.\\
		Xét dãy $1^2;\, 2^2;\, 3^2;\, 4^2;  \cdots \Rightarrow \dfrac{u_2}{u_1}=4 \neq \dfrac{9}{4}=\dfrac{u_3}{u_2}$ nên không phải là cấp số nhân.
		\begin{itemize}
			\item $2;\, 4;\, 8;\, 16; \ldots$ là cấp số nhân và $u_n=2^n$. 
			\item $1;\,-1;\,1;\,-1; \ldots$ là  cấp số nhân và $u_n=(-1)^n$.
			\item Ta có $\dfrac{2^2}{1^2}\neq \dfrac{3^2}{2^2}$ nên dãy số $1^2;\, 2^2;\, 3^2;\, 4^2; \ldots$ không phải là cấp  số nhân.
			\item $a;\,  a^3;\, a^5;\, a^7;\, \ldots(a\ne 0)$ là cấp số nhân và $u_n=a^{2 n-1}=\dfrac{1}{a} \cdot\left(a^2\right)^n$.
		\end{itemize}
	} 
\end{ex}

% nguồn: Ok_TracNghiemLop11_Lan1
\begin{ex}%[1D2N3-2]%[ma:f3f22501]
	Trong các dãy số dưới đây,  dãy số nào \textbf{không} là cấp số nhân?
	 \choice
		{$2;0;0;0;0;\ldots$} 
		{$1;3;9;27;81;\ldots$} 
		{\True $1;2;3;4;5;\ldots$} 
		 {$2;4;8;16;32;\ldots$}
	\loigiai{ 
		Dãy số $2;0;0;0;0;\ldots$ có $u_1=2,  q=0$.\\
		 Dãy số $1;3;9;27;81;\ldots$ có $u_1=1, q=3$.\\
		Dãy số $1;2;3;4;5;\ldots$ có $\dfrac{u_2}{u_1}=\dfrac{2}{1}\neq\dfrac{u_3}{u_2}=\dfrac{3}{2}$. Đây không là  cấp số nhân.\\
		Dãy só $2;4;8;16;32;\ldots$  có $u_1=2, q=2$.
	} 
\end{ex}

% nguồn: Ok_TracNghiemLop11_Lan1
\begin{ex}%[1D2N3-2]%[ma:f273f22e]
	Cho cấp  số nhân $\left(u_n\right)$ với $u_1=3$ và $u_2=15$. Công bội của cấp số nhân đã cho bằng
	\choice 
		{\True  $ 5$}
		{$ 12$} 
		{$\dfrac{1}{5}$} 
		 {$ 4$}
	\loigiai{ 
	 Ta có $u_2=u_1\cdot q\Rightarrow q=\dfrac{u_2}{u_1}=\dfrac{15}{3}=5$.}
\end{ex}

% nguồn: Ok_TracNghiemLop11_Lan1
\begin{ex}%[1D2N3-2]%[ma:dea9c920]
	Trong các dãy số hữu hạn sau, dãy số nào là cấp số  nhân?
	 \choice
		{$2; 4; 8; 16; 32;  63$}
		 {\True $1;-2;4;-8;16;-32$}
		{$1; 3;  9; 27; 54; 162$}
		{$4; 2;  1; \dfrac{1}{2}; \dfrac{1}{4}; \dfrac{1}{16}$}
	\loigiai{ 
		Vì  $\dfrac{4}{2}= \dfrac{8}{4} = \dfrac{16}{8}= \dfrac{32}{16} \neq \dfrac{63}{32}$ nên dãy số $2; 4; 8; 16; 32; 63$ không là cấp số nhân.\\
		Vì $\dfrac{3}{1}= \dfrac{9}{3} = \dfrac{27}{9} \ne \dfrac{54}{27}$ nên dãy số $1; 3; 9; 27; 54; 162$ không là cấp  số nhân.\\
		Vì $\dfrac{2}{4}= \dfrac{1}{2} = \dfrac{\dfrac{1}{2}}{1} = \dfrac{\dfrac{1}{4}}{\dfrac{1}{2}} \ne  \dfrac{\dfrac{1}{16}}{\dfrac{1}{4}}$ nên dãy số ở $4; 2; 1; \dfrac{1}{2}; \dfrac{1}{4}; \dfrac{1}{16}$ không là cấp số nhân.\\
		Vì $\dfrac{-2}{1}= \dfrac{4}{-2} = \dfrac{-8}{4} = \dfrac{16}{-8} = \dfrac{-32}{16}=-2$ nên dãy số $1;-2;4;-8;16;-32$ là  cấp số nhân.
	} 
\end{ex}

% nguồn: Ok_TracNghiemLop11_Lan1
\begin{ex}%[1D2N3-2]%[ma:7cac6643]
	Cho  ba mệnh đề
	\begin{itemize}
		\item $M$:  Dãy số $2$, $6$, $18$, $54$ là một cấp số nhân.
		\item $N$: Dãy số $48$, $24$, $12$, $6$,  $3$ là một cấp số nhân.
		\item $Q$: Dãy  số $1$, $-1$, $1$, $-1$, $1$, $-1$ là một cấp số nhân.
	\end{itemize}
	 Trong ba mệnh đề trên, có bao nhiêu mệnh đề đúng?
	\choice 
		{$1$} 
		 {\True $3$}
		 {$0$}
		{$2$} 
	\loigiai{ 
		Ta  có
		\begin{itemize}
			\item $M$: Dãy số $2$, $6$, $18$, $54$ là một  cấp số nhân có số hạng đầu bằng $u_1=2$, công bội $q=3$.
			\item $N$: Dãy số $48$, $24$, $12$, $6$, $3$  là một cấp số nhân có số hạng đầu bằng $u_1=48$, công bội $q=\dfrac{1}{2}$.
			\item $Q$: Dãy số $1$, $-1$, $1$, $-1$, $1$, $-1$ là một cấp số nhân  có số hạng đầu bằng $u_1=1$, công bội $q=-1$.
		\end{itemize}
		Do  đó có $3$ mệnh đều đúng.
	} 
\end{ex}

% nguồn: Ok_TracNghiemLop11_Lan1
\begin{ex}%[1D2N3-2]%[ma:7f630dd1]
	Cho cấp số nhân có số hạng đầu tiên là $u_1 = 9$, $q = 2$. Số hạng thứ $2$ bằng 
	 \choice
		 {$90$}
		{\True  $18$}
		{$36$} 
		{$16$} 
	\loigiai{ 
		Ta có $$u_2 = u_1  \cdot q^1 = 9 \cdot 2^1 = 18.$$
	} 
\end{ex}

% nguồn: Ok_TracNghiemLop11_Lan1
\begin{ex}%[1D2N3-2]%[ma:3ffcdd49]
	Dãy số nào sau đây {\bf không} phải là  cấp số nhân?
	 \choice
		 {$4$;$2$;$\dfrac{1}{2}$;$\dfrac{1}{4}$;$\ldots $}
		{\True $\dfrac{1}{\pi}$;$\dfrac{1}{{\pi}^2}$;$\dfrac{1}{{\pi}^4}$;$\dfrac{1}{{\pi  }^6}$;$\ldots$}
		 {$1$;$2$;$4$;$8$;$\ldots$}
		{$3$;$3^2$;$3^3$;$3^4$;$\ldots$} 
	 \loigiai{
		Xét $\dfrac{1}{\pi}$;$\dfrac{1}{{\pi}^2}$;$\dfrac{1}{{\pi}^4}$;$\dfrac{1}{{\pi  }^6}$;$\ldots$.\\
		Ta thấy  $\dfrac{u_2}{u_1}=\dfrac{1}{\pi}\ne \dfrac{1}{{\pi }^2}=\dfrac{u_3}{u_2}$.
	 }
\end{ex}

% nguồn: Ok_TracNghiemLop11_Lan1
\begin{ex}%[1D2N3-2]%[ma:92bfb064]
	Dãy số nào sau đây {\bf không} phải là  một cấp số nhân?
	\choice 
		{$1$; $-3$; $9$;  $-27$; $54$}
		{$1$; $2$; $4$; $8$;  $16$}
		{\True $1$; $3$; $5$; $7$;  $9$}
		{$1$; $-2$; $4$; $-8$; $16$} 
	\loigiai{ 
		\begin{itemize}
			\item $1$; $-3$; $9$; $-27$; $54$  là cấp số nhân với công bội $q=-3$.
			\item $1$; $2$; $4$; $8$; $16$ là cấp số  nhân với công bội $q=2$.
			\item $1$; $-2$; $4$; $-8$; $16$  là cấp số nhân với công bội $q=-2$.
			\item $1$; $3$; $5$; $7$; $9$ là cấp  số cộng với công sai $d=2$.
		\end{itemize}
	} 
\end{ex}

% nguồn: Ok_TracNghiemLop11_Lan1
\begin{ex}%[1D2N3-2]%[ma:8a3a2989]
	Trong các dãy số sau, dãy số nào là cấp số  nhân?
	\choice 
		{\True $128$; $-64$; $32$; $-16$;  $8$}
		 {$\sqrt{2}$; $2$; $4$; $4\sqrt{2}$}
		{$5$; $6$; $7$; $8$} 
		{$15$;  $5$; $1$; $\dfrac{1}{5}$}
	 \loigiai{
		Ta thấy $\dfrac{128}{-64}=\dfrac{-64}{32}=\dfrac{32}{-16}=\dfrac{-16}{8}=-2$ nên dãy số $128$; $-64$; $32$; $-16$; $8$ là một cấp số nhân với  công bội bằng $-2$.
	} 
\end{ex}

% nguồn: Ok_TracNghiemLop11_Lan1
\begin{ex}%[1D2N3-2]%[ma:fab87e97]
	Cho cấp số nhân $32$; $16$; $8$; $4$; $2$ có công bội là 
	 \choice
		 {$q=2$}
		{\True  $q=\dfrac{1}{2}$}
		 {$q=\dfrac{1}{4}$}
		 {$q=\dfrac{1}{3}$}
	 \loigiai{
		Cho cấp số nhân $32$; $16$; $8$;  $4$; $2$ có công bội là $q=\dfrac{1}{2}$.
	} 
\end{ex}

% nguồn: Ok_TracNghiemLop11_Lan1
\begin{ex}%[1D2N3-2]%[ma:e02fdcd5]
	 Dãy số $1;\, 2;\, 4;\, 8;\, 16;\, 32; \ldots$ là một cấp số nhân với
	 \choice
		{Công bội là $3$  và số hạng đầu tiên là $1$}
		{\True Công bội là $2$ và số hạng đầu tiên là  $1$}
		{Công bội là $4$ và số hạng  đầu tiên là $2$}
		{Công bội là $2$ và số hạng đầu tiên là $2$} 
	\loigiai{ 
		Cấp số nhân $1;\, 2;\, 4;\, 8;\, 16;\, 32; \ldots$ có $\heva{& u_1=1 \\& q=\dfrac{u_2}{u_1}=2}$. 
	} 
\end{ex}

% nguồn: Ok_TracNghiemLop11_Lan1
\begin{ex}%[1D2N3-2]%[ma:909e1261]
	 Dãy số nào sau đây \textbf{không} phải là cấp số nhân?
	\choice 
		{$1;\, 2;\, 4;\,  8; \ldots$}
		 {$3;\, 3^2;\, 3^3;\, 3^4;\, \ldots$}
		{$4;\, 2;\,  \dfrac{1}{2};\, \dfrac{1}{4}; \ldots$}
		{\True $\dfrac{1}{\pi};\, \dfrac{1}{\pi^2};\,  \dfrac{1}{\pi^4};\, \dfrac{1}{\pi^6}; \ldots$}
	 \loigiai{
		Xét $1;\, 2;\, 4;\, 8; \ldots$ là cấp số  nhân công bội là $2$.\\
		Xét $3;\, 3^2;\, 3^3;\, 3^4;\, \ldots$ là cấp số nhân công  bội là $3$.\\
		Xét $4;\, 2;\, \dfrac{1}{2};\, \dfrac{1}{4}; \ldots$ là cấp số nhân công bội là $\dfrac{1}{2}$.\\ 
		Xét $\dfrac{1}{\pi};\, \dfrac{1}{\pi^2};\, \dfrac{1}{\pi^4};\, \dfrac{1}{\pi^6}; \ldots$ không là cấp số nhân vì $\dfrac{u_2}{u_1}=\dfrac{1}{\pi} \neq  \dfrac{1}{\pi^2}=\dfrac{u_3}{u_2}$.\\
	 }
\end{ex}

% nguồn: Ok_TracNghiemLop11_Lan1
\begin{ex}%[1D2N3-2]%[ma:ef45f7f1]
	Trong các dãy số sau đây dãy số nào là cấp số nhân? 
	 \choice
		{\True  Dãy số $-2,\,2,\,-2,\,2, \ldots,-2,2,-2,2, \ldots$}
		{Dãy số $\left(u_n\right)$, xác định bởi công thức $u_n=3^n+1$ với $n \in \mathbb{N^*}$} 
		{Dãy số $\left(u_n\right)$, xác  định bởi hệ: $\heva{&u_1=1 \\ &u_n=u_{n-1}+2\left(n \in \mathbb{N^*};\, n \geq 2\right)}$}
		{Dãy số các số tự nhiên  $1,\, 2,\, 3,\ldots$}
	\loigiai{ 
		Dãy số $-2,\,2,\,-2,\,2, \ldots,-2,2,-2,2, \ldots$ là cấp số  nhân với số hạng đầu $u_1=-2$ , công bội $q=-1$.\\
		Dãy số $\left(u_n\right)$  xác định bởi công thức $u_n=3^n+1$ có $u_1=3^1+1=4$, $u_2=3^2+1=10$, $u_3=3^3+1=28$.\\
		\begin{nx}
			$\dfrac{u_3}{u_2} \neq \dfrac{u_2}{u_1}$ nên $\left(u_n\right)$ không là  cấp số nhân.
		\end{nx}
		Dãy số $\left(u_n\right)$,  xác định bởi hệ $\heva{&u_1=1 \\ &u_n=u_{n-1}+2\left(n \in \mathbb{N*};\, n \geq 2\right)}$ có $u_1=1, u_1=3, u_3=5$.\\
		\begin{nx}
			$\dfrac{u_3}{u_2} \neq \dfrac{u_2}{u_1}$  nên $\left(u_n\right)$ không là cấp số nhân.
		\end{nx}
		Dãy  số các số tự nhiên $1,2,3, \ldots$. có $u_1=1, u_1=2, u_3=3$.
		\begin{nx}
			$\dfrac{u_3}{u_2} \neq \dfrac{u_2}{u_1}$ nên không  là cấp số nhân.
		\end{nx}
	 }
\end{ex}

% nguồn: Ok_TracNghiemLop11_Lan1
\begin{ex}%[1D2N3-2]%[ma:1e245bc1]
	Trong các dãy số $(u_n)$ sau  đây, dãy số nào là cấp số nhân?
	\choice 
		{$u_n=(-1)^n\cdot  n$}
		{$u_n=n^2$} 
		{\True $u_n=2^n$} 
		{$u_n=\dfrac{n}{3^n}$} 
	 \loigiai{
		Lập  tỉ số $\dfrac{u_{n+1}}{u_n}$.
		\begin{itemize}
			\item $\dfrac{u_{n+1}}{u_n}=\dfrac{(-1)^{n+1}\cdot (n+1)}{(-1)^n\cdot n}=-\dfrac{n+1}{n}\Rightarrow (u_n)$ không phải cấp số nhân.\\ 
			 \item $\dfrac{u_{n+1}}{u_n}=\dfrac{(n+1)^2}{n^2}\Rightarrow (u_n)$ không phải là cấp số nhân.\\
			\item $\dfrac{u_{n+1}}{u_n}=\dfrac{2^{n+1}}{2^n}=2\Rightarrow u_{n+1}=2u_n\Rightarrow (u_n)$ là cấp số nhân có công bội bằng $2$.\\ 
			\item $\dfrac{u_{n+1}}{u_n}=\dfrac{n+1}{3n}\Rightarrow (u_n)$ không phải là  cấp số nhân.
		\end{itemize}
	 }
\end{ex}

% nguồn: Ok_TracNghiemLop11_Lan1
\begin{ex}%[1D2N3-2]%[ma:237d29fa]
	Trong  các dãy số sau, dãy số nào là cấp số nhân?
	 \choice
		{$1$, $2$, $4$,  $6$}
		{$1$, $3$,  $12$, $60$}
		 {\True $-1$, $4$, $-16$, $64$}
		{$-1$,  $-5$, $-25$, $125$}
	 \loigiai{
		Dãy $-1$, $4$, $-16$, $64$ là cấp số nhân với  số hạng đầu là $-1$, công bội $q=-4$.
	} 
\end{ex}

% nguồn: Ok_Toan11_TracNghiem_Lan2
\begin{ex}%[1D2H3-2]%[ma:f7f2c472]
	Dãy số nào sau đây \textbf{không} phải là cấp số nhân? 
	\choice 
		{\True $1$; $-3$; $9$;  $-27$; $54$}
		{$1$;  $2$; $4$; $8$; $16$}
		{$1$; $-1$;  $1$; $-1$; $1$}
		{$1$;  $-2$; $4$; $-8$; $16$}
	 \loigiai{
		\begin{itemize}
			\item Dãy $1$;  $-3$; $9$; $-27$; $54$ là không là cấp số nhân vì
			\[-3=1\cdot (-3);~(-27)\cdot (-3)=81\neq 54.\] 
			\item Dãy  $1$; $2$; $4$; $8$; $16$ là cấp số nhân với $q=2$.
			\item Dãy $1$; $-1$; $1$; $-1$; $1$ là cấp số  nhân với $q=-1$.
			\item Dãy $1$; $-2$;  $4$; $-8$; $16$ là cấp số nhân với $q=-2$.
		\end{itemize}
	} 
\end{ex}

% nguồn: Ok_Toan11_TracNghiem_Lan2
\begin{ex}%[1D2H3-2]%[ma:c7c4b565]
	Cho dãy số $(u_n)$ là một cấp số nhân với $u_n\neq 0$, $n\in \mathbb{N}^*$.  Dãy số nào sau đây \textbf{không phải} là cấp số nhân?
	\choice 
		{$u_1; u_3; u_5; \ldots$} 
		{$3u_1; 3u_2; 3u_3;  \ldots$}
		{$\dfrac{1}{u_1}; \dfrac{1}{u_2}; \dfrac{1}{u_3}; \ldots$} 
		{\True  $u_1+2; u_2+2; u_3+2; \ldots$}
	\loigiai{ 
		Giả sử  $(u_n)$ là cấp số nhân công bội $q$, thì $u_{n+1}=q\cdot u_n$, $\forall n\in \mathbb{N}^*$.
		\begin{itemize}
			\item Dãy $u_1; u_3; u_5; \ldots$ là cấp  số nhân công bội $q^2$.
			\item Dãy $3u_1; 3u_2;  3u_3; \ldots$ là cấp số nhân công bội $2q$.
			\item Dãy $\dfrac{1}{u_1}; \dfrac{1}{u_2}; \dfrac{1}{u_3}; \ldots$ là cấp số nhân công bội $\dfrac{1}{q}$. 
			\item Dãy $u_1+2; u_2+2; u_3+2; \ldots$  không phải là cấp số nhân.
		\end{itemize}
	 }
\end{ex}

% nguồn: Ok_Toan11_TracNghiem_Lan2
\begin{ex}%[1D2H3-2]%[ma:c90985d0]
	Cho dãy số $\left(u_n\right)$ là một cấp số nhân với $u_n \ne 0$, $n \in \mathbb{N}^*$. Dãy số nào sau đây \textbf{không} phải là  cấp số nhân?
	\choice 
		{$u_1$; $u_3$; $u_5$; $\ldots$} 
		{$3u_1$; $3u_2$; $3u_3$;  $\ldots$}
		{$\dfrac{1}{u_1}$;  $\dfrac{1}{u_2}$; $\dfrac{1}{u_3}$; $\ldots$}
		{\True $u_1 + 2$; $u_2 + 2$; $u_3  + 2$; $\ldots$}
	 \loigiai{
		Giả sử $\left(u_n\right)$  là cấp số nhân công bội $q$, thì
		\begin{itemize}
			\item Dãy $u_1$; $u_3$; $u_5$; $\ldots$  là cấp số nhân công bội $q^2$.
			\item Dãy $3u_1$; $3u_2$; $3u_3$; $\ldots$ là cấp số nhân công bội $q$. 
			\item Dãy $\dfrac{1}{u_1}$; $\dfrac{1}{u_2}$; $\dfrac{1}{u_3}$; $\ldots$ là cấp số nhân công bội $\dfrac{1}{q}$. 
			\item Dãy $u_1 + 2$; $u_2 +  2$; $u_3 + 2$; $\ldots$ không phải là cấp số nhân.
		\end{itemize}
	 }
\end{ex}

% nguồn: Ok_Toan11_TracNghiem_Lan2
\begin{ex}%[1D2H3-2]%[ma:696af2f3]
	Trong các số sau,  dãy số nào là một cấp số nhân
	 \choice
		{\True $1$; $-3$;  $9$; $-27$; $81$}
		 {$1$; $-3$; $-6$; $-9$; $-12$}
		{$1$; $-2$; $-4$; $-8$;  $-16$}
		{$0$;  $3$; $9$; $27$; $81$}
	\loigiai{ 
		Dãy số $1$; $-3$; $9$; $-27$; $81$ là cấp số nhân vì 
		 có công bộ $q = -3$.
	 }
\end{ex}

% nguồn: Ok_Toan11_TracNghiem_Lan2
\begin{ex}%[1D2H3-2]%[ma:a66aa99e]
	Trong các  dãy số sau, dãy nào là cấp số nhân?
	\choice 
		{$u_n=(-1)^n n$} 
		{$u_n=n^2$} 
		{\True $u_n=2^n$} 
		{$u_n=\dfrac{n}{3^n}$} 
	 \loigiai{
		Lập tỉ số $\dfrac{u_{n+1}}{u_n}$\\ 
		\begin{itemize}
			\item $\dfrac{u_{n+1}}{u_n}=\dfrac{(-1)^{n+1}\cdot \left(n+1\right)}{(-1)^n\cdot  n}=-\dfrac{n+1}{n} \Rightarrow (u_n)$ không phải cấp số nhân.
			\item $\dfrac{u_{n+1}}{u_n}=\dfrac{(n+1)^2}{n^2} \Rightarrow (u_n)$ không phải là cấp  số nhân.
			\item $\dfrac{u_{n+1}}{u_n}=\dfrac{2^{n+1}}{2^n}=2\Rightarrow{u_{n+1}}=2u_n \Rightarrow (u_n)$ là cấp  số nhân có công bội bằng $2$.
			\item $\dfrac{u_{n+1}}{u_n}=\dfrac{n+1}{3n} \Rightarrow  (u_n)$ không phải là cấp số nhân.
	\end{itemize}}
\end{ex}

% nguồn: Ok_Toan11_TracNghiem_Lan2
\begin{ex}%[1D2H3-2]%[ma:906f81ba]
	Trong các dãy số $(u_n)$ sau đây, dãy số nào là cấp  số nhân?
	\choice 
		 {$u_n=3n$}
		 {\True $u_n=2^n$}
		 {$u_n=\dfrac{1}{n}$}
		 {$u_n=2^n+1$}
	 \loigiai{
		Ta thấy,  với $\forall n\ge 2$, $n\in \mathbb{N}^*$ dãy số $(u_n)=2^n$ có tính chất: $\dfrac{u_n}{u_{n-1}}=\dfrac{2^n}{2^{n-1}}=2$ nên là cấp số nhân với công bội $q=2$, $u_1=2$.
	 }
\end{ex}

% nguồn: Ok_Toan11_TracNghiem_Lan2
\begin{ex}%[1D2H3-2]%[ma:f9e6ebee]
	Dãy số nào sau đây \textbf{không} phải là cấp số nhân? 
	 \choice
		{$1$; $-1$; $1$;  $-1$}
		{\True $1$;  $-3$; $9$; $10$}
		{$1$; $0$; $0$;  $0$}
		{$32$;  $16$; $8$; $4$}
	 \loigiai{
		 Nếu $\left(u_n\right)$ là cấp số nhân với công bội $q$ ta có
		\[u_{n+1}=u_n\cdot q  \Rightarrow q=\dfrac{u_{n+1}}{u_n}.\]
		\begin{itemize}
			\item $1$; $-1$; $1$; $-1$ là cấp số nhân  với $q=-1$.
			\item $1$; $-3$; $9$;  $10$ không là cấp số nhân.
			\item $1$; $0$; $0$; $0$ là cấp số nhân với  $q=0$.
			 \item $32$; $16$; $8$; $4$ là cấp số nhân với $q=\dfrac{1}{2}$.
		\end{itemize}
	} 
\end{ex}

% nguồn: Ok_Toan11_TracNghiem_Lan2
\begin{ex}%[1D2H3-2]%[ma:10027863]
	Trong dãy số $\left(u_n\right)$ sau, dãy nào là  cấp số nhân?
	 \choice
		{$u_n  = n^2 + n + 1$}
		{$u_n = (n + 2)  \cdot 3^n$}
		 {$\heva{
		&u_1 =  2\\
		&u_{n+1} = \dfrac{6}{u_n},\ \forall n \in \mathbb{N}^*}$ 
		 }
		{\True  $u_n = (-4)^{2n+1}$}
	\loigiai{ 
		\begin{itemize}
			\item Với $u_n = n^2 +  n + 1$\\
			$\Rightarrow \dfrac{u_{n+1}}{u_n} = \dfrac{n^2  + 3n + 3}{n^2 + n + 1}$, không phải là hằng số nên không là cấp số nhân.
			\item Với $u_n  = (n + 2)\cdot 3^n$\\
			$\Rightarrow \dfrac{u_{n+1}}{u_n} = \dfrac{(n + 3)\cdot 3^{n+1}}{(n +  2)\cdot 3^n} = \dfrac{3(n + 3)}{n + 2}$, cũng không là hằng số nên không là cấp số nhân.
			\item Với dãy số  cho bởi
			 $\heva{
			&u_1 =  2\\
			 &u_{n+1} = \dfrac{6}{u_n}}$
			$\Rightarrow u_2 = \dfrac{6}{2} = 3$; $u_3 = \dfrac{6}{3} = 2$; $u_4 = \dfrac{6}{2} = 3$; $\ldots$\\ 
			Vì $\dfrac{u_3}{u_2}\neq \dfrac{u_2}{u_1}$. 
			Nên  $\left(u_n\right)$ không phải cấp số nhân.
			\item  Với $u_n = (-4)^{2n+1} \Rightarrow \dfrac{u_{n+1}}{u_n} = \dfrac{(-4)^{2(n+1)+1}}{(-4)^{2n+1}} = (-4)^2 = 16$.\\
			Vậy dãy là cấp số nhân với công bội $q = 16$. 
		\end{itemize}
	 }
\end{ex}

% nguồn: Ok_Toan11_TracNghiem_Lan2
\begin{ex}%[1D2H3-2]%[ma:16742b8b]
	Chọn  cấp số nhân trong các dãy số sau
	 \choice
		{$1$;  $0{,}2$; $0{,}4$; $0{,}008$; $\ldots$}
		 {$2$; $22$; $222$; $2222$; $\ldots$}
		{$x$; $2x$; $3x$;  $4x$; $\ldots$}
		 {\True $1$; $-x^2$; $x^4$; $-x^6$; $\ldots$}
	\loigiai{ 
		Dãy  số: $1;\ -x^2;\ x^4;\ -x^6;\ldots$ là cấp số nhân với
		\begin{itemize}
			\item Số hạng đầu $u_1 = 1$ 
			\item Công bội  $q = \dfrac{-x^2}{1} = -x^2$
		\end{itemize}
	} 
\end{ex}

% nguồn: Ok_Toan11_TracNghiem_Lan2
\begin{ex}%[1D2H3-2]%[ma:fe7722b3]
	Dãy số  nào sau đây \textbf{không} phải là cấp số nhân?
	\choice 
		{$1$;  $2$; $4$; $8$; $\ldots$}
		{$3$; $3^2$; $3^3$; $3^4$; $\ldots$} 
		{$4$; $2$; $\dfrac{1}{2}$;  $\dfrac{1}{4}$; $\ldots$}
		{\True $\dfrac{1}{\pi}$;  $\dfrac{1}{\pi^2}$; $\dfrac{1}{\pi^4}$; $\dfrac{1}{\pi^6}$; $\ldots$}
	 \loigiai{
		 Xét đáp án $\mathrm{D}$, ta có
		 \[\dfrac{u_2}{u_1}=\dfrac{\dfrac{1}{\pi^2}}{\dfrac{1}{\pi}}=\dfrac{1}{\pi}\neq \dfrac{1}{\pi^2}=\dfrac{\dfrac{1}{\pi^4}}{\dfrac{1}{\pi^2}}=\dfrac{u_3}{u_2}.\]
		Vậy dãy $\dfrac{1}{\pi}$;  $\dfrac{1}{\pi^2}$; $\dfrac{1}{\pi^4}$; $\dfrac{1}{\pi^6}$; $\ldots$ \textbf{không} phải là cấp số nhân.
	 }
\end{ex}

% nguồn: Ok_Toan11_TracNghiem_Lan2
\begin{ex}%[1D2H3-2]%[ma:821b1083]
	Trong các dãy số sau,  dãy số nào là một cấp số nhân?
	\choice 
		{$1$; $2$; $3$; $4$; $5$; $6$;  $\ldots$}
		{$2$; $4$,$6$; $8$;  $16$; $32$; $\ldots$}
		{$-2$;  $-3$,$-4$; $-5$; $-6$; $-7$; $\ldots$}
		 {\True $1$; $2$,$4$; $8$; $16$; $32$; $\ldots$}
	 \loigiai{
		Nhận thấy $\dfrac{u_2}{u_1}\neq \dfrac{u_3}{u_2}$ nên các dãy số ở đáp án $\mathrm{A}$, $\mathrm{B}$ và $\mathrm{C}$  không phải cấp số nhân.\\
		Riêng đối với dãy $1$; $2$,$4$; $8$; $16$; $32$;  $\ldots$ ở đáp án $\mathrm{D}$ thỏa mãn $u_{n+1}=2\cdot u_n$, $\forall n \in \mathbb{N}^*$.\\
		Vậy dãy  số $1$; $2$; $4$; $8$; $16$; $32$; $\ldots$ là cấp số nhân với $u_1=1$ và công bội $q=2$.
	} 
\end{ex}

% nguồn: Ok_Toan11_TracNghiem_Lan2
\begin{ex}%[1D2H3-2]%[ma:96c1fc28]
	Dãy  nào sau đây là một cấp số nhân?
	\choice 
		{$1$, $2$, $3$, $4$,  $\ldots$}
		{$1$, $3$, $5$,  $7$, $\ldots$}
		{\True $2$, $4$, $8$, $16$,  $\ldots$}
		{$2$, $4$,  $6$, $8$, $\ldots$}
	 \loigiai{
		Ta có $2$, $4$,  $8$, $16$, $\ldots$ là cấp số nhân có số hạng đầu $u_1 = 2$ và công bội $q = 2$.
	} 
\end{ex}

% nguồn: Ok_Toan11_TracNghiem_Lan2
\begin{ex}%[1D2H3-2]%[ma:8e2d7d45]
	Trong các dãy số đã cho dưới đây, dãy số nào là  cấp số nhân?
	\choice 
		 {$1$; $2$; $3$; $4$; $5$}
		 {$1$; $3$; $6$; $9$; $12$}
		{$2$; $4$; $6$;  $8$; $10$}
		{\True $2$; $2$; $2$; $2$;  $2$}
	\loigiai{Ta thấy ở đáp án $\mathrm{D}$ có $u_1=u_2=u_3=u_4=u_5=2$ nên đây là cấp số nhân với  công bội $q=1$.
	} 
\end{ex}

% nguồn: Ok_Toan11_TracNghiem_Lan2
\begin{ex}%[1D2H3-2]%[ma:7d49917b]
	Cho cấp số nhân có các số hạng lần  lượt là $\dfrac{1}{4}$; $\dfrac{1}{2}$; $1$; $\ldots$; $2048$. Tính tổng $S$ của tất cả các số hạng của cấp số nhân đã cho.
	 \choice
		{$S =  2\,047{,}75$}
		 {$S = 2\,049{,}75$}
		 {\True $S = 4\,095{,}75$}
		{$S  = 4\,096{,}75$}
	 \loigiai{
		Cấp số nhân đã cho  có $\heva{&u_1 = \dfrac{1}{4}\\ &q = 2} \Rightarrow 2048 = 2^{11} = u_1 q^{n - 1} = \dfrac{1}{4} \cdot 2^{n - 1} = 2^{n - 3} \Leftrightarrow n = 14$.\\
		Do đó cấp  số nhân đã cho có tất cả $14$ số hạng.\\
		 Vậy $S_{14} = u_1 \cdot \dfrac{1 - q^{14}}{1 - q} = \dfrac{1}{4} \cdot \dfrac{1 - 2^{14 }}{1 - 2} = 4\,095{,}75$.
	 }
\end{ex}

% nguồn: Ok_Toan11_TracNghiem_Lan2
\begin{ex}%[1D2H3-2]%[ma:d34a03f1]
	Cho dãy số $(u_n)$ có số hạng tổng quát là $u_n=3\cdot 2^{n+1}$, $n \in \mathbf{N}^{*}$. Chọn kết  luận đúng
	 \choice
		 {\True Dãy số là cấp số nhân có số hạng đầu $u_1=12$}
		 {Dãy số là cấp số cộng có công sai $d=2$}
		{Dãy số là cấp số cộng có số hạng đầu $u_1=6$} 
		{Dãy số là cấp số nhân  có công bội $q=3$}
	\loigiai{ 
		Dãy số $(u_n)$ có số  hạng tổng quát là $u_n=3\cdot 2^{n+1}, n \in \mathbf{N}^{*} \Rightarrow u_{n+1}=3 \cdot 2^{n+2}$.\\
	Xét thương $\dfrac{u_{n+1}}{u_n}=\dfrac{3\cdot 2^{n+2}}{3\cdot 2^{n+1}}=2$ với $n \in \mathbf{N}^{*}$ nên dãy số $(u_n)$ là một cấp số  nhân có công bội $q=2$ và có số hạng đầu là $u_1=3\cdot 2^{1+1}=12$.}
\end{ex}

% nguồn: Ok_Toan11_TracNghiem_Lan2
\begin{ex}%[1D2H3-2]%[ma:27ede16c]
	Cho dãy số $(u_n)$ xác định bởi $\heva{  & u_1=-2 \\ &u_{n+1}=-\dfrac{1}{10}u_n }$. Chọn hệ thức đúng trong các hệ thức sau
	\choice 
		{\True $(u_n)$ là cấp  số nhân có công bội $q=-\dfrac{1}{10}$}
		 {$u_n=(-2)\dfrac{1}{10^{n-1}}$}
		{$u_n=\dfrac{u_{n-1}+u_{n+1}}{2}$ (với  $n\ge 2$)}
		{$u_n=\sqrt{u\dfrac{n-1}\cdot u_{n+1}{}}$  (với $n\ge 2$)}
	\loigiai{ 
	Ta có  $\dfrac{u_{n+1}}{u_n}=-\dfrac{1}{10}$ nên $(u_n)$ là cấp số nhân có công bội $q=-\dfrac{1}{10}$.
	 }
\end{ex}

% nguồn: Ok_Toan11_TracNghiem_Lan2
\begin{ex}%[1D2H3-2]%[ma:8f5ade93]
	Trong các dãy số $(u_n)$ cho bởi số hạng tổng quát $u_n$ sau, dãy số nào  là một cấp số nhân?
	 \choice
		{\True $u_n = \dfrac{1}{3^{n - 2}}$} 
		 {$u_n = \dfrac{1}{3^{n - 1}}$}
		{$u_n = n  + \dfrac{1}{3}$}
		 {$u_n = n^2 - \dfrac{1}{3}$}
	\loigiai{ 
		Ta có $u_n  = \dfrac{1}{3^{n - 2}} = 9 \cdot \left(\dfrac{1}{3}\right)^n$
		 là cấp số nhân với $u_1 = 3$ và công bội $q = \dfrac{1}{3}$.
	} 
\end{ex}

% nguồn: Ok_Toan11_TracNghiem_Lan2
\begin{ex}%[1D2H3-2]%[ma:f3622370]
	 Cho cấp số nhân có $u_2=\dfrac{1}{4}$, $u_5=16$. Tìm số hạng đầu $u_1$ và công bội $q$.
	\choice 
		{$u_1=\dfrac{1}{2}$, $q=\dfrac{1}{2}$} 
		 {$u_1=-\dfrac{1}{2}$, $q=-\dfrac{1}{2}$}
		{\True $u_1=\dfrac{1}{16}$, $q=4$} 
		 {$u_1=-\dfrac{1}{16}$, $q=-4$}
	\loigiai{ 
	Ta có $\heva{ & u_1q=\dfrac{1}{4} \\ &u_1q^4=16 }\Rightarrow  q^3=64\Rightarrow q=4$.\\
	Vậy  $u_1=\dfrac{1}{16}$.
	 }
\end{ex}

% nguồn: Ok_Toan11_TracNghiem_Lan2
\begin{ex}%[1D2H3-2]%[ma:c5df5867]
	Trong các dãy số sau,  dãy số nào \textbf{không} phải là một cấp số nhân?
	\choice 
		{$2$; $4$;  $8$; $16$; $\ldots$}
		{$1$; $-1$; $1$;  $-1$; $\ldots$}
		{\True $1^2$; $2^2$;  $3^2$; $4^2$; $\ldots$}
		 {$a$; $a^3$; $a^5$; $a^7$; $\ldots \left(a\neq 0\right)$}
	 \loigiai{
		Xét đáp án $1^2$; $2^2$;  $3^2$; $4^2$; $\ldots$\\
		Suy ra  \[\dfrac{u_2}{u_1}=4\neq \dfrac{9}{4}=\dfrac{u_3}{u_2}.\]
		Vậy dãy số $1^2$; $2^2$; $3^2$; $4^2$; $\ldots$ \textbf{không} phải là cấp số  nhân.
	 }
\end{ex}

% nguồn: Ok_Toan11_TracNghiem_Lan2
\begin{ex}%[1D2H3-2]%[ma:4ebd82ed]
	Trong các dãy số $(u_n)$ sau đây, dãy số nào là cấp số  nhân?
	\choice 
		{$u_n=(-1)^nn$} 
		 {$u_n=n^2$}
		 {\True $u_n=2^n$}
		 {$u_n=\dfrac{n}{3^n}$}
	 \loigiai{
		Lập tỉ số $\dfrac{u_{n+1}}{u_n}$  ta có
		\begin{itemize}
			\item Dãy  số $u_n=(-1)^nn$ có $\dfrac{u_{n+1}}{u_n}=\dfrac{(-1)^{n+1}\cdot (n+1)}{(-1)^n\cdot n}=-\dfrac{n+1}{n} \Rightarrow (u_n)$ không phải cấp số nhân.
			\item Dãy số $u_n=n^2$ có $\dfrac{u_{n+1}}{u_n}=\dfrac{(n+1)^2}{n^2} \Rightarrow (u_n)$  không phải là cấp số nhân.
			\item Dãy số $u_n=2^n$ có  $\dfrac{u_{n+1}}{u_n}=\dfrac{2^{n+1}}{2^n}=2 \Rightarrow (u_n)$ là cấp số nhân có công bội bằng $2$.
			\item Dãy số  $u_n=\dfrac{n}{3^n}$ có $\dfrac{u_{n+1}}{u_n}=\dfrac{n+1}{3n} \Rightarrow (u_n)$ không phải là cấp số nhân.
		\end{itemize}
	} 
\end{ex}

% nguồn: Ok_Toan11_TracNghiem_Lan2
\begin{ex}%[1D2H3-2]%[ma:8eb0fbf4]
	Trong các dãy số sau, dãy nào  là cấp số nhân?
	 \choice
		 {$u_n = (-1)^n \cdot n$}
		{$u_n =  n^2$}
		{\True  $u_n = 2^n$}
		 {$u_n = \dfrac{n}{3^n}$}
	\loigiai{ 
		Lập tỉ số $\dfrac{u_{n+1}}{u_n}$. 
		\begin{itemize}
			\item Với  $u_n = (-1)^n \cdot n$ \\
			\[\dfrac{u_{n+1}}{u_n} = \dfrac{(-1)^{n+1} \cdot (n+1)}{(-1)^n \cdot  n} = \dfrac{- (n+1)}{n}\]
			$\Rightarrow \left(u_n\right)$ không phải cấp  số nhân.
			 \item Với $u_n = n^2$ \\
			$\dfrac{u_{n+1}}{u_n}  = \dfrac{(n+1)^2}{n^2} \Rightarrow \left(u_n\right)$ không phải là cấp số nhân.
			\item Với $u_n  = 2^n$ \\
			$\dfrac{u_{n+1}}{u_n} = \dfrac{2^{n+1}}{2^n} = 2  \Rightarrow u_{n+1} = 2u_n \Rightarrow \left(u_n\right)$ là cấp số nhân có công bội bằng $2$.
			\item Với $u_n =  \dfrac{n}{3^n}$: \\
			$\dfrac{u_{n+1}}{u_n} = \dfrac{n+1}{3n}  \Rightarrow \left(u_n\right)$ không phải là cấp số nhân.
		\end{itemize}
	} 
\end{ex}

% nguồn: Ok_Toan11_TracNghiem_Lan2
\begin{ex}%[1D2H3-2]%[ma:447aeb06]
	Cho dãy số: $-1$; $\dfrac{1}{3};\ -\dfrac{1}{9}$; $\dfrac{1}{27}$;  $-\dfrac{1}{81}$; $\ldots$. Khẳng định nào sau đây là \textbf{sai}?
	 \choice
		{\True Dãy số không phải là một  cấp số nhân}
		{Dãy số này là  cấp số nhân có $u_1 = -1$; $q = -\dfrac{1}{3}$}
		{Số hạng tổng quát $u_n  = (-1)^n \cdot \dfrac{1}{3^{n-1}}$}
		{Là dãy số không tăng, không  giảm}
	\loigiai{ 
		 Ta có $\dfrac{1}{3}=-1\cdot \left(-\dfrac{1}{3}\right)$; $-\dfrac{1}{9}=-\dfrac{1}{3}\cdot \left(-\dfrac{1}{3}\right)$;$\dfrac{1}{27}=-\dfrac{1}{9}\cdot \left(-\dfrac{1}{3}\right)$; $\ldots \ldots$\\
		Vậy  dãy trên là cấp số nhân với $u_1=-1$; $q=-\dfrac{1}{3}$.\\
		Áp dụng công thức số  hạng tổng quát cấp số nhân ta có
		\[u_n=u_11^{n-1}=-1\left(-\dfrac{1}{3}\right)^{n-1}=\left(-1\right)^n\cdot \dfrac{1}{3^{n-1}}.\] 
	 }
\end{ex}

% nguồn: Ok_Toan11_TracNghiem_Lan2
\begin{ex}%[1D2H3-2]%[ma:82e1e5f0]
	Số hạng đầu và công bội $q$  của cấp số nhân với $u_7=-5$, $u_{10}=135$ là
	\choice 
		 {$u_1=\dfrac{5}{729}$, $q=-3$}
		{$u_1=-\dfrac{5}{729}$,  $q=3$}
		 {$u_1=\dfrac{5}{729}$, $q=3$}
		 {\True $u_1=-\dfrac{5}{729}$, $q=-3$}
	\loigiai{ 
		Vì $\left(u_n\right)$ là cấp số nhân nên $u_7=u_1\cdot q^6=-5$,  $u_{10}=u_1\cdot q^9=135$.\\
		 Suy ra $\dfrac{u_{10}}{u_7}=\dfrac{135}{-5}\Leftrightarrow\dfrac{u_1q^9}{u_1q^6}=-27 \Rightarrow q=-3$.\\
		Khi  đó $u_1=\dfrac{u_7}{q^6}=-\dfrac{5}{729}$.
	} 
\end{ex}

% nguồn: Ok_Toan11_TracNghiem_Lan2
\begin{ex}%[1D2H3-2]%[ma:6f68bcc9]
	Trong các dãy  số $(u_n)$ sau, dãy nào là cấp số nhân?
	\choice 
		 {$u_n=n^2+n+1$}
		 {$u_n=(n+2)\cdot 3^n$}
		 {$\heva{&u_1=2 \\&u_{n+1}=\dfrac{6}{u_n},\, \forall n\in \mathbb{N}^*}$}
		{\True  $u_n=(-4)^{2n+1}$}
	\loigiai{ 
		\begin{itemize}
			\item Dãy số $u_n=n^2+n+1$ có $\dfrac{u_{n+1}}{u_n}=\dfrac{n^2+3n+3}{n^2+n+1}$,  $\forall n\in \mathbb{N}^*$, không phải là hằng số. Vậy $(u_n)$ không phải là cấp số nhân.
			\item Dãy số $u_n=(n+2)\cdot 3^n$ có $\dfrac{u_{n+1}}{u_n}=\dfrac{(n+3)\cdot 3^{n+1}}{(n+2)\cdot 3^n}=\dfrac{3(n+3)}{n+2}$, $\forall n\in  \mathbb{N}^*$, không phải là hằng số. Vậy $(u_n)$ không phải là cấp số nhân.
			\item Dãy số $\heva{&u_1=2 \\&u_{n+1}=\dfrac{6}{u_n},\, \forall n\in \mathbb{N}^*}$  có $u_1=2; u_2=3; u_3=2; u_4=3; \ldots$ \\
			Vì $\dfrac{u_3}{u_2}\neq \dfrac{u_2}{u_1}$ nên $(u_n)$ không phải là  cấp số nhân.
			\item Dãy số $u_n=(-4)^{2n+1}$  có $\dfrac{u_{n+1}}{u_n}=\dfrac{(-4)^{2(n+1)+1}}{(-4)^{2n+1}}=16$, $\forall n\in \mathbb{N}^*$. Vậy $(u_n)$ là một cấp số nhân.
		\end{itemize}
	 }
\end{ex}

% nguồn: Ok_Toan11_TracNghiem_Lan2
\begin{ex}%[1D2H3-2]%[ma:66e9e816]
	Cho cấp số nhân $(u_n)$ có tổng $n$ số hạng đầu tiên là $S_n=5^n-1$ với $n=1{,}2,\cdots$. Tìm số hạng đầu $u_1$  và công bội $q$ của cấp số nhân đó?
	 \choice
		 {$u_1=5$, $q=4$}
		 {$u_1=5$, $q=6$}
		{\True $u_1=4$,  $q=5$}
		{$u_1=6$, $q=5$} 
	\loigiai{Ta  có
		\[\heva{&u_1=S_1=5^1-1=4\\&u_1+u_2=S_2=5^2-1=24} \Leftrightarrow  \heva{&u_1=4\\&u_2=24-u_1=20} \Rightarrow u_1=4, q=\dfrac{u_2}{u_1}=5.\]
	 }
\end{ex}

% nguồn: Ok_TracNghiemLop11_Lan1
\begin{ex}%[1D2H3-2]%[ma:9bfd6aa3]
	Cho cấp số nhân $\left(u_n\right)$ thỏa mãn  $u_1=3,\, u_5=48$. Công bội của cấp số nhân bằng
	 \choice
		{$16$} 
		 {$-2$}
		 {$2$}
		 {\True$\pm 2$}
	 \loigiai{
		Gọi $q$ là công bội  của cấp số nhân $\left(u_n\right)$.\\
		Với $u_1=3$, $u_5=48$ suy ra $\heva{&u_1=3 \\& u_1 \cdot  q^4=48} \Leftrightarrow\heva{&u_1=3 \\& q^4=16} \Leftrightarrow\heva{&u_1=3 \\ &q= \pm 2}$.\\
		Vậy công bội của  cấp số nhân $\left(u_n\right)$ là $q= \pm 2$.
	 }
\end{ex}

% nguồn: Ok_TracNghiemLop11_Lan1
\begin{ex}%[1D2H3-2]%[ma:42b05f64]
	Trong các dãy số sau, dãy số  nào là một cấp số nhân?
	\choice 
		 {$1, 2, 3, 4, 5, 6, \ldots$}
		{$2, 4, 6,  8, 16, 32,\ldots$}
		{$-2, -3, -4, -5, -6, -7,  \ldots$}
		{\True $1, 2, 4,  8, 16, 32, \ldots$}
	 \loigiai{
		Nhận thấy dãy $1, 2, 4, 8, 16,  32, \ldots$ thỏa mãn: $u_{n+1}=2.u_n\, \,\forall n\in {{\mathbb{N}}^*}$.\\
		Vậy dãy số $1, 2, 4, 8, 16, 32, \ldots$ là  cấp số nhân với $u_1=1$ và công bội $q=2$.
	 }
\end{ex}

% nguồn: Ok_TracNghiemLop11_Lan1
\begin{ex}%[1D2H3-2]%[ma:4cbebc16]
	Cho cấp số nhân $\left(u_n\right)$ có $u_2=6$  và $u_3=18$. Tìm công bội của cấp số nhân đã cho.
	\choice 
		 {$12$}
		{\True  $3$}
		{$-12$} 
		 {$24$}
	\loigiai{ 
		Do $u_2$ và $u_3$ là hai số hạng liên tiếp của cấp số nhân nên  suy ra công bội $q=\dfrac{u_3}{u_2}=\dfrac{18}{6}=
		 3$.
	} 
\end{ex}

% nguồn: Ok_TracNghiemLop11_Lan1
\begin{ex}%[1D2H3-2]%[ma:5f114f09]
	Cho cấp số nhân $\heva{&u_1+u_5=51\\&u_2+u_6=102}$. Số hạng đầu $u_1$ bằng  bao nhiêu?
	 \choice
		{$u_1  =2$}
		{\True $u_1 =3$} 
		 {$u_1 =6$}
		{$u_1  =4$}
	 \loigiai{
		Ta có  số hạng tổng quát của cấp số nhân là $u_n = u_1 \cdot q^{n-1}$. Khi đó
		\[\heva{&u_1+u_5=51\\&u_2+u_6=102} \Leftrightarrow \heva{&u_1+u_1 q^4=51\\&u_1q+u_1 q^4=102} \Leftrightarrow \heva{&u_1 \left(1+q^4\right)=51\\&u_1q\left(1+q^4\right)=102}  \Leftrightarrow \heva{&u_1 =3\\&q=2.}\]
		Vậy số hạng đầu $u_1 =3$ và công bội  $q=2$.
	 }
\end{ex}

% nguồn: Ok_TracNghiemLop11_Lan1
\begin{ex}%[1D2H3-2]%[ma:84575ad0]
	Người ta trồng $2047$ cây theo hình tam giác như sau: hàng thứ nhất có $1$ cây,  hàng thứ hai có $2$ cây, hàng thứ ba có $4$ cây,\ldots , hàng thứ $n$ có $2^{n-1}$ cây. Người ta có thể trồng bao nhiêu hàng cây với cách trồng này?
	\choice 
		{\True $11$  hàng}
		{$12$  hàng}
		 {$10$ hàng}
		 {$15$ hàng}
	 \loigiai{
	Theo  đề bài ta có
	\begin{eqnarray*}
		&&1+2+4+\ldots +2^{n-1}=2047\\ 
		&\Leftrightarrow  &2^{n}-1=2047\\
		&\Leftrightarrow &n=11. 
	\end{eqnarray*}
	Vậy người ta  có thể trồng $11$ hàng cây với cách trồng này.
	} 
\end{ex}

% nguồn: Ok_TracNghiemLop11_Lan1
\begin{ex}%[1D2H3-2]%[ma:02b79a3f]
	Cho cấp số  nhân $(u_n)$ với $u_1 = 3$; $u_2 = 9$. Công bội của cấp số nhân đã cho bằng
	\choice 
		{$6$} 
		 {\True $3$}
		{$12$} 
		 {$-6$}
	 \loigiai{
		Công bội của cấp số  nhân $q=\dfrac{u_{n+1}}{u_n}$=$\dfrac{9}{3}=3$.
	 }
\end{ex}

% nguồn: Ok_TracNghiemLop11_Lan1
\begin{ex}%[1D2H3-2]%[ma:6e9be4b8]
	Dãy số cho bởi công thức nào dưới đây \textbf{không} phải là cấp  số nhân?
	\choice 
		{$u_n=\dfrac{3^n}{2}$} 
		{$u_n=\dfrac{2}{5^n}$} 
		 {$u_n=\left(-1\right)^n$}
		{\True  $u_n=3n+2$}
	\loigiai{ 
		Ta có 
		$u_n=\dfrac{3^n}{2}$  là số hạng tổng quát của cấp số nhân vì $\dfrac{u_{n+1}}{u_n}=3$.\\
		$u_n=\dfrac{2}{5^n}$ là số hạng tổng quát của cấp số nhân  vì $\dfrac{u_{n+1}}{u_n}=\dfrac{1}{5}$.\\
		$u_n=\left(-1\right)^n$ là số hạng tổng quát của cấp  số nhân vì $\dfrac{u_{n+1}}{u_n}=-1$.
	 }
\end{ex}

% nguồn: Ok_TracNghiemLop11_Lan1
\begin{ex}%[1D2H3-2]%[ma:649f1e02]
	Cho cấp số nhân $(u_{n})$ có $u_{7}=-5$ và  $u_{10}=135$. Công bội của cấp số nhân là
	\choice{\True  $q=-3$}
		 {$q=-\dfrac{1}{3}$}
		{$q=3$} 
		{$q=9$} 
	 \loigiai{
		Ta có $\heva{& u_7=u_1\cdot q^6=-5\\& u_{10}=u_1\cdot q^9=135}\Leftrightarrow\heva{& u_1=-\dfrac{5}{729}\\&  q=-3.}$
	} 
\end{ex}

% nguồn: Ok_TracNghiemLop11_Lan1
\begin{ex}%[1D2H3-2]%[ma:e4bb6063]
	Trong các dãy số sau, đây số nào là cấp  số nhân?
	 \choice
		{\True  $u_n=\dfrac{1}{3^{n-2}}$}
		 {$u_n=\dfrac{1}{3^n}-1$}
		{$u_n=n+\dfrac{1}{3}$} 
		 {$u_n=n^2+\dfrac{1}{3}$}
	 \loigiai{
		Dãy số $u_n=\dfrac{1}{3^{n-2}}$ là cấp số  nhân với số hạng đầu $u_1=3$, công bội $q=\dfrac{1}{3}$.
	 }
\end{ex}

% nguồn: Ok_TracNghiemLop11_Lan1
\begin{ex}%[1D2H3-2]%[ma:668e9030]
	Cho dãy số $(u_n)$ có số  hạng tổng quát là $u_n=3\cdot 2^{n+1}\; \left(\forall n\in {{\mathbb{N}}^*}\right)$. Chọn kết luận đúng:
	 \choice
		{\True Dãy số là cấp số nhân có số hạng đầu  $u_1=12$}
		{Dãy số là cấp số cộng có  công sai $d=2$}
		{Dãy số  là cấp số cộng có số hạng đầu $u_1=6$}
		{Dãy số  là cấp số nhân có công bội $q=3$}
	 \loigiai{
		Dãy số $(u_n)$có số hạng tổng quát  là $u_n=3\cdot 2^{n+1}\;\left(\forall n\in \mathbb{N}^*\right)\Rightarrow u_{n+1}=3\cdot 2^{n+2}$.\\
		Xét thương $\dfrac{u_{n+1}}{u_n}=\dfrac{3\cdot 2^{n+2}}{3\cdot 2^{n+1}}=2=const$ với $\forall n\in {{\mathbb{N}}^*}$ nên dãy số $(u_n)$ là một cấp số nhân có công bội $q=2$ và có số hạng đầu  là $u_1=3\cdot 2^{1+1}=12$.
	} 
\end{ex}

% nguồn: Ok_TracNghiemLop11_Lan1
\begin{ex}%[1D2H3-2]%[ma:aaccea49]
	Cho cấp số nhân  $(u_n)$ có $u_1=\dfrac{1}{2}$, $u_2=16$. Khi đó công bội $q$ là
	 \choice
		 {$64$}
		 {$8$}
		 {$4$}
		{\True  $32$}
	 \loigiai{
		Cấp số nhân $(u_n)$ có công bội là $q=\dfrac{u_2}{u_1}=\dfrac{16}{\dfrac{1}{2}}=32$. 
	} 
\end{ex}

% nguồn: Ok_TracNghiemLop11_Lan1
\begin{ex}%[1D2H3-2]%[ma:bfd73e23]
	Cho cấp số nhân $(u_n)$ có $u_2=-6$ và $u_6=-486$. Tìm công bội  $q$ của cấp số nhân đã cho, biết rằng $u_3>0$.
	 \choice
		{$q=-3$} 
		 {$q=-\dfrac{1}{3}$}
		{$q=\dfrac{1}{3}$} 
		{\True  $q=3$}
	 \loigiai{
		Ta  có
		$$\heva{& -6=u_2=u_1 q  \\ & -486=u_6=u_1 q^5=u_1 q \cdot q^4=-6\cdot q^4 }\Rightarrow q^4=81=3^4 \quad q=3.$$
	 }
\end{ex}

% nguồn: Ok_Toan11_TracNghiem_Lan2
\begin{ex}%[1D2V3-2]%[ma:4073d9c0]
	Dãy số nào sau  đây là cấp số nhân?
	\choice 
		{$ 
		\heva{ 
		 &u_1 = 1\\
		&u_{n+1} = u_n + 1,\ n  \ge 1
		 }
		$} 
		{\True $ 
		\heva{ 
		 &u_1 = -1\\
		&u_{n+1}  = -3u_n,\ n \ge 1
		 }
		 $}
		{$ 
		\heva{ 
		&u_1 =  -2\\
		&u_{n+1} = 2u_n + 3,\ n  \ge 1
		} 
		 $}
		{$ 
		\heva{ 
		&u_1 =  \dfrac{\pi}{2}\\
		&u_n = \sin\left(\dfrac{\pi}{n -  1}\right),\ n \ge 1
		} 
		$} 
	\loigiai{ 
		Ta có $u_{n+1} = -3u_n$ là cấp số nhân với công bội $q = -3$. 
	 }
\end{ex}

% nguồn: Ok_Toan11_TracNghiem_Lan2
\begin{ex}%[1D2V3-2]%[ma:e7398bc0]
	Dãy  số nào sau đây là cấp số nhân?
	 \choice
		 {$\heva{&u_1=1 \\&u_{n+1}=u_n+1,\, n\ge 1}$}
		{\True $\heva{&u_1=-1 \\&u_{n+1}=-3u_n,\,  n\ge 1}$}
		{$\heva{&u_1=-2 \\&u_{n+1}=2u_n+3,\, n\ge  1}$}
		{$\heva{&u_1=\dfrac{\pi}{2}  \\ &u_n=\sin \left(\dfrac{\pi}{n-1}\right),\, n\ge 1}$}
	 \loigiai{
		Dãy số $(u_n)$ là cấp  số nhân $\Leftrightarrow u_{n+1}=qu_n$. \\
		Vậy dãy số $\heva{&u_1=-1 \\&u_{n+1}=-3u_n,\, n\ge 1}$ là  CSN.
	 }
\end{ex}

% nguồn: Ok_Toan11_TracNghiem_Lan2
\begin{ex}%[1D2V3-2]%[ma:10511475]
	Cho dãy số $\left(u_n\right)$ được xác định bởi $\heva{&u_1=2\\&u_{n+1}=4u_n+9}$, $\forall n\ge 1$, dãy số $\left(v_n\right)$ xác định bởi $v_n=u_n+3$, $\forall n\ge 1$ là một cấp số  nhân. Số hạng đầu và công bội của cấp số nhân đó là
	\choice 
		 {$v_1=3$, $q=5$}
		 {$v_1=4$, $q=5$}
		 {$v_1=3$, $q=4$}
		{\True $v_1=5$,  $q=4$}
	 \loigiai{
		 Vì có $v_n=u_n+3\quad(1)\Rightarrow v_{n+1}=u_{n+1}+3\quad(2)$.\\
		Theo  đề $u_{n+1}=4u_n+9\Rightarrow u_{n+1}+3=4\left(u_n+3\right)$.\\
		Thay $(1)$ và $(2)$ vào được $v_{n+1}=4v_n$, $\forall n\ge 1\Rightarrow \dfrac{v_{n+1}}{v_n}=4$.\\ 
		Kết luận $\left(v_n\right)$ là  cấp số nhân với công bội $q=4$ và số hạng đầu $v_1=u_1+3=5$.
	} 
\end{ex}


%%%%%%%%%%%%%%%% ĐÚNG/SAI %%%%%%%%%%%%%%%%

% nguồn: Ok_dungsai_lop11_Lan1
\begin{ex}%[1D2N3-2]%[ma:16072c3b]
 Cho hình chóp $S.ABCD$ có đáy $ABCD$ là hình  bình hành, điểm $M$ di động trên cạnh $AD$. Một mặt phẳng $(\alpha)$ qua $M$ và song song với hai đường thẳng $CD$, $SA$, cắt $BC$, $SC$ và $SD$ lần lượt tại $N$, $P$, $Q$. Khi đó
  \choiceTF
  {Giao tuyến của mặt phẳng $(\alpha)$ với mặt phẳng $(ABCD)$ là đường thẳng  đi qua $M$ và song song với $AD$}
  {\True Giao tuyến của mặt phẳng  $(\alpha)$ với mặt phẳng $(SAD)$ là đường thẳng đi qua $M$ và song song với $SA$}
  {\True Tứ giác $MNPQ$ là hình thang có hai đáy là $MN$ và $P Q$} 
  {Gọi $I=MQ \cap NP$. Khi đó $I$ thuộc đường thẳng đi qua  $S$ và song song với $AB$}
  \loigiai{
  \begin{center}
   \begin{tikzpicture}[scale=0.7,font=\footnotesize,line join=round,line cap=round,>=stealth]
      \path
       (0:0) coordinate (A)
      (0:5) coordinate (D) 
      (-135:2.8)  coordinate (B)
      +(D) coordinate (C) 
      (80:5) coordinate  (S)
      ($(B)!0.6!(C)$)  coordinate (N)
      ($(A)!0.6!(D)$) coordinate (M) 
       ($(S)!0.6!(D)$) coordinate (Q)
      ($(S)!0.6!(C)$) coordinate  (P)
       (intersection of N--P and M--Q) coordinate (I)
       ($(S)!2!(I)$) coordinate (T) node[above] {$x$}
      (intersection of I--P and S--Q) coordinate (T1) 
      ($(S)!0.9!(T1)$) coordinate  (T2);
     \draw[dashed] (B)--(A)--(D) 
      (S)--(A) 
       (N)--(M)--(Q)--(T2);
      \draw[thick] (S)--(B)
       (S)--(C)
      (S)--(T2) 
       (Q)--(D)
      (C)--(B) 
      (D)--(C) 
       (N)--(P)--(Q)
      (S)--(T) 
       (Q)--(I)--(P);
     \foreach \x/\g in {A/135,B/180,C/-30,D/0,S/90,M/-37,N/-90,Q/10,P/180,I/90}\draw[fill=white] (\x)  circle (.03) +(\g:.3) node{$\x$};
    \end{tikzpicture}
  \end{center}
  \begin{itemchoice}
   \itemch Sai. Ta có $\heva{&M  N=(\alpha) \cap(A B C D) \\ &C D \parallel(\alpha) \\ &C D \subset(A B C D)} \Rightarrow(\alpha) \cap(A B C D)=M N \parallel C D\quad(1)$ .
   \itemch Đúng. Ta có $\heva{&M Q=(\alpha) \cap(S A D) \\ &S A \parallel(\alpha) \\ &S A \subset(S A D)} \Rightarrow(\alpha)  \cap(SAD)=MQ \parallel SA$.
   \itemch  Đúng. Ta có $
   \heva{&PQ=(\alpha)  \cap(SCD) \\
   &CD  \parallel(\alpha) \\
   &CD \subset(S C D)} \Rightarrow(\alpha) \cap(S C  D)=P Q \parallel C D\quad(2)$.
    \\Từ (1) và (2) suy ra tứ giác $MNPQ$ là hình thang có hai đáy là $MN$ và $PQ$.
   \itemch Sai. Xét $(SAD) \cap(SBC)$. 
    Ta có $\heva{&S \in(SAD) \cap(SBC) \\ &AD \parallel BC \\ &AD \subset(SAD), BC \subset(S B C)} \Rightarrow(S A D) \cap(S B C)=S x$
   (với $S x$ qua $S$ và  $S x \parallel AD \parallel BC$).\\
   Vì $\heva{ &I \in N P, N P \subset(S  B C) \\
   &I \in M Q, M Q \subset(S A D)}\Rightarrow I \in(S  A D) \cap(S B C)$.
    Suy ra $I \in S x$ (với $S x$ cố định)..
  \end{itemchoice}
  }
\end{ex}

% nguồn: Ok_dungsai_lop11_Lan1
\begin{ex}%[1D2N3-2]%[ma:dad9f8fa]
 Hãy cho biết tính đúng sai của mỗi mệnh đề  dưới đây?
 \choiceTF 
  {\True Ba số $2$; $10$; $50$ theo thứ tự lập  thành một cấp số nhân}
  {Dãy  số $0,\ -3,\ -9,\ -27,\ -81$ là một cấp số nhân.}
  {\True Dãy số $u_n=3^n$ là một cấp  số nhân}
  {Dãy số $u_n=n^2-\dfrac{1}{4}$ là một cấp số  nhân}
 \loigiai{ 
  \begin{itemchoice}
   \itemch Ba số $2$; $10$; $50$ theo thứ tự lập thành cấp số nhân có  công bội $ q=5 $.
   \itemch Dãy số $0,\ -3,\ -9,\ -27,\ -81$ không là  một cấp số nhân.
   \itemch Xét dãy $(u_n)$ với $u_n=3^n\Rightarrow u_{n+1}=3^{n+1}=3\cdot 3^n=3u_n$. Vậy  $u_n=3^n$ là cấp số nhân.
   \itemch $\dfrac{u_{n+1}}{u_n}=\dfrac{(n+1)^2-\dfrac{1}{4}}{n^2-\dfrac{1}{4}}=\dfrac{(n+1)^2-\dfrac{1}{4}}{n^2-\dfrac{1}{4}}=\dfrac{4n^2+8n+3}{4n^2-1},\forall n\in \mathbb{N^*}$, không  là hằng số.\\
   Vậy $(u_n)$ không là  cấp số nhân.
  \end{itemchoice}
 } 
\end{ex}

% nguồn: Ok_dungsai_lop11_Lan1
\begin{ex}%[1D2N3-2]%[ma:8810bcb9]
 Cho dãy số $\left(u_n\right)$, biết $\heva{&u_1=-1\\&u_{n+1}=u_n+3}$  với $n \ge 1$. Khi đó
 \choiceTF 
   {\True Bốn số hạng đầu tiên của dãy số lần lượt là $-1$; $2$; $5$; $8$}
  {Số hạng thứ năm  của dãy là $13$}
  {Công thức số hạng tổng  quát của dãy số là $u_n=2 n-3$}
  {\True $101$ là số hạng thứ  $35$ của dãy số đã cho}
  \loigiai{
  \begin{itemchoice}
   \itemch Đúng. Ta có  $u_1=-1$; $u_2=u_1+3=2$; $u_3=u_2+3=5$; $u_4=u_3+3=8$.
    \itemch Sai. $u_5=u_4+3=11$.
   \itemch Sai. Từ giả thiết, ta có $\heva{&u_1=-1\\&u_2=u_1+3\\&u_3=u_2+3\\ &\cdots \cdots \cdots \\  &u_n=u_{n-1}+3}$.\\
   Cộng theo vế toàn bộ các đẳng thức trên và triệt tiêu các số hạng giống nhau  ở hai vế, ta có $u_n=-1+3(n-1)=3 n-4$.\\Vậy công thức số hạng tổng quát của dãy số là $u_n=3 n-4$.
   \itemch Đúng. Xét $101=3 n-4 \Rightarrow n=35$.\\ Vậy $101$ là số hạng  thứ $35$ của dãy số đã cho.
  \end{itemchoice}
 } 
\end{ex}

% nguồn: Ok_dungsai_lop11_Lan1
\begin{ex}%[1D2N3-2]%[ma:a17384de]
 Cho phương trình lượng giác $2 \sin \left(x-\dfrac{\pi}{12}\right)+\sqrt{3}=0$. Khi  đó
 \choiceTF 
  {Phương trình tương  đương $\sin \left(x-\dfrac{\pi}{12}\right)=\sin \left(\dfrac{\pi}{3}\right)$}
  {Phương trình có nghiệm là: $x=\dfrac{\pi}{4}+k 2 \pi; x=\dfrac{7 \pi}{12}+k 2 \pi(k \in \mathbb{Z})$} 
  {\True Phương trình có nghiệm âm lớn nhất  bằng $-\dfrac{\pi}{4}$}
  {\True  Số nghiệm của phương trình trong khoảng $(-\pi, \pi)$ là hai nghiệm}
  \loigiai{
  \begin{itemchoice}
   \itemch Sai. Ta có 
   \begin{eqnarray*}
     & & 2 \sin \left(x-\frac{\pi}{12}\right)+\sqrt{3}=0\\
    &\Leftrightarrow &  \sin \left(x-\dfrac{\pi}{12}\right)=-\dfrac{\sqrt{3}}{2}\\
    &\Leftrightarrow  & \sin \left(x-\dfrac{\pi}{12}\right)=\sin \left(-\dfrac{\pi}{3}\right)
   \end{eqnarray*}.
    \itemch Sai.
    $
   \Leftrightarrow\left[\begin{array}{ l }
   { x - \dfrac { \pi }{ 1 2 } = - \dfrac { \pi }{ 3 } + k 2  \pi } \\
   { x - \dfrac { \pi }{ 1 2 } = \pi - ( - \dfrac { \pi }{  3 }) + k 2 \pi }
  \end{array} \quad ( k \in \mathbb { Z }) \Leftrightarrow \left[\begin{array}{l}
  x=-\dfrac{\pi}{4}+k 2 \pi  \\
  x=\dfrac{17 \pi}{12}+k 2 \pi 
 \end{array}(k \in \mathbb{Z})\right.\right.
  $.
 Vậy  phương trình có nghiệm là $x=-\dfrac{\pi}{4}+k 2 \pi$ và $x=\dfrac{17 \pi}{12}+k 2 \pi(k \in \mathbb{Z})$.
 \itemch Đúng. Phương trình có nghiệm âm lớn  nhất bằng $-\dfrac{\pi}{4}$.
 \itemch Đúng. Số nghiệm của phương trình  trong khoảng $(-\pi, \pi)$ là hai nghiệm.
 \end{itemchoice}
} 
\end{ex}

% nguồn: Ok_dungsai_lop11_Lan1
\begin{ex}%[1D2N3-2]%[ma:58f471bb]
  Cho các dãy số sau đây ${u_n=(\sqrt{5})^{2 n-3}}$; ${v_n=\dfrac{2}{n}}$;${w_n=\dfrac{3^{n+1}}{2^n}}$ và dãy số hữu hạn gồm các số hạng ${16; 4; 1; \dfrac{1}{4}; \dfrac{1}{16}; \dfrac{1}{64}}$. Khi đó
  \choiceTF
   {\True ${\left(u_n\right)}$ là một cấp số nhân công bội ${q=5}$}
  {\True ${\left(v_n\right)}$ không phải là một cấp số  nhân}
  {\True ${\left(w_n\right)}$ là một cấp số  nhân có số hạng đầu $w=\dfrac{9}{2}$}
  {Dãy số hữu hạn đã cho theo thứ tự lập thành cấp số nhân có  công bội bằng $\dfrac{1}{8}$}
 \loigiai{ 
  \begin{enumerate}
   \item Ta có $u_{n+1}={{(\sqrt{5})}^{2(n+1)-3}}={{(\sqrt{5})}^{2n-1}}\Rightarrow  \dfrac{{{u}_{n+1}}}{{{u}_{n}}}=\dfrac{{{(\sqrt{5})}^{2n-1}}}{{{(\sqrt{5})}^{2n-3}}}={{(\sqrt{5})}^{2}}=5,\forall n\in \mathbb{N}^{*}$.\\
   Do đó ${\left(u_n\right)}$ là một cấp số nhân  có số hạng đầu ${u_1=\dfrac{1}{\sqrt{5}}}$ với công bội ${q=5}$.
   \item  Ta có ${v_{n+1}=\dfrac{2}{n+1} \Rightarrow \dfrac{v_{n+1}}{v_n}=\dfrac{n}{n+1}}$ (tỉ số này còn phụ thuộc vào ${n}$).\\
   Do đó  ${\left(v_n\right)}$ không phải là một cấp số nhân.
   \item Ta có  $w_{n+1}=\dfrac{3^{n+2}}{2^{n+1}}=\dfrac{3}{2} \cdot \dfrac{3^{n+1}}{2^n} \Rightarrow \dfrac{w_{n+1}}{w_n}=\dfrac{\dfrac{3}{2} \cdot \dfrac{3^{n+1}}{2^n}}{\dfrac{3^{n+1}}{2^n}}=\dfrac{3}{2}, \forall n \in \mathbb{N}^*$.\\
   Do đó  ${\left(w_n\right)}$ là một cấp số nhân có số hạng đầu ${w_1=\dfrac{3^2}{2^1}=\dfrac{9}{2}}$ với công bội ${q=\dfrac{3}{2}}$
   \item Đặt $u_1=16; u_2=4; u_3=1;  u_4=\dfrac{1}{4}; u_5=\dfrac{1}{16}; u_6=\dfrac{1}{64}$.\\
   Ta có ${u_2=u_1 \cdot \dfrac{1}{4};\, u_3=u_2 \cdot \dfrac{1}{4};\, u_4=u_3 \cdot \dfrac{1}{4};\, u_5=u_4 \cdot \dfrac{1}{4};\, u_6=u_5 \cdot \dfrac{1}{4}}$.  \\
    Vì vậy dãy số hữu hạn đã cho theo thứ tự lập thành cấp số nhân có công bội bằng ${\dfrac{1}{4}}$.
  \end{enumerate}
 } 
\end{ex}

% nguồn: Ok_dungsai_lop11_Lan1
\begin{ex}%[1D2N3-2]%[ma:8a6d3fd9]
  Cho góc lượng giác $x$, biết $\tan x=\dfrac{1}{3}$ với $\dfrac{\pi}{2}<x<\pi$. Khi đó
 \choiceTF 
   {\True $\cos x<0$}
  {$\cos x=-\dfrac{\sqrt{10}}{10}$} 
  {\True  $\sin x=-\dfrac{\sqrt{10}}{10}$}
  {$\sin  x+\cos x=-\dfrac{\sqrt{10}}{5}$}
  \loigiai{
  \begin{itemchoice}
   \itemch Đúng. Ta có $\tan x=\frac{1}{3} \Rightarrow \cot x=\dfrac{1}{\tan x}=3$. Vì  $\dfrac{\pi}{2}<x<\pi$ nên $\cos x<0$.
   \itemch Sai. Ta có $\dfrac{1}{\cos ^2 x}=1+\tan ^2 x=1+\dfrac{1}{9}=\dfrac{10}{9} \Rightarrow \cos ^2 x=\dfrac{9}{10}  \Rightarrow \cos x=-\dfrac{3 \sqrt{10}}{10}$.
   \itemch  Đúng. $\sin x=\cos x \cdot\tan x=-\dfrac{3 \sqrt{10}}{10} \cdot \dfrac{1}{3}=-\dfrac{\sqrt{10}}{10}$ .
   \itemch Sai. $\sin  x+\cos x=-\dfrac{2 \sqrt{10}}{5}$.
  \end{itemchoice}
 } 
\end{ex}

% nguồn: Ok_dungsai_lop11_Lan1
\begin{ex}%[1D2N3-2]%[ma:d6b981c5]
 Hãy cho biết tính đúng  sai của mỗi mệnh đề dưới đây?
 \choiceTF 
  {\True  Cấp số nhân với $u_1=2$; $u_2=8$. Giá trị của công bội $q$ bằng $4$}
  {\True Cấp  số nhân $u_{n}=3^{n-1}$ có công bội $q=3$}
  {\True Cấp số  nhân $\left(u_{n}\right)$ có số hạng đầu $u_{1}=2$ và $u_{6}=486$ có công bội $q=3$}
  {Một cấp số nhân có số $ u_3=2 $ và $ u_6=54 $ có công bội $  q=-3 $}
 \loigiai{ 
  \begin{itemchoice}
   \itemch Ta có  $u_2=u_1\cdot q$. Suy ra $q=\dfrac{u_2}{u_1}=\dfrac{8}{2}=4$.\\
   Vậy công bội $q=4$. 
   \itemch Ta thấy $u_{n+1} =3 u_n$ cho nên công bội của cấp số nhân là $q=3$. 
   \itemch Ta có $\heva{&u_1=2\\&u_6=486} \Leftrightarrow \heva{&u_1=2\\&u_1  \cdot q^5=486}$ $\Rightarrow q^5=243\Rightarrow q=3$.
    \itemch Ta có $ \heva{&u_3=2\\&u_6=54}\Leftrightarrow \heva{&u_1\cdot q^2=2\\&u_1 \cdot q^5=54} \Rightarrow q^3=27 \Leftrightarrow q=3$.
  \end{itemchoice}
  }
\end{ex}

% nguồn: Ok_dungsai_lop11_Lan1
\begin{ex}%[1D2H3-2]%[ma:276138da]
 Cho cấp số nhân ${\left(u_n\right)}$  với công bội ${q<0}$ và ${u_2=4, u_4=9}$. Khi đó
 \choiceTF 
  {\True Số hạng  đầu ${u_1=-\dfrac{8}{3}}$}
  {Số hạng  $u_5=\dfrac{27}{2}$}
  {$-\dfrac{2187}{32}$ là số  hạng thứ $8$}
  {Cấp số nhân có công bội ${q=-\dfrac{3}{2}}$} 
 \loigiai{ 
  Ta có ${u_2=u_1 q=4,\,  u_4=u_1 q^3=9 \Rightarrow \dfrac{u_4}{u_2}=\dfrac{u_1 q^3}{u_1 q} \Rightarrow \dfrac{9}{4}=q^2 \Rightarrow q=-\dfrac{3}{2}(q<0)}$.\\
  Thay ${q=-\dfrac{3}{2}}$ vào ${u_2}$, ta được  ${u_1\left(-\dfrac{3}{2}\right)=4 \Rightarrow u_1=-\dfrac{8}{3}}$.\\
  Vậy cấp số nhân đã cho có số hạng đầu ${u_1=-\dfrac{8}{3}}$ và công bội  ${q=-\dfrac{3}{2}}$.\\
   Khi đó ${{u}_{n}}=-\dfrac{8}{3}\cdot{{\left( -\dfrac{3}{2} \right)}^{n-1}}$.\\
  Vậy  $u_5=-\dfrac{27}{2}$.\\
  Do $-\dfrac{2187}{32}\ne -\dfrac{8}{3}\cdot{{\left( -\dfrac{3}{2} \right)}^{7}}$ nên không  phải là số hạng thứ $8$.
  }
\end{ex}

% nguồn: Ok_dungsai_lop11_Lan1
\begin{ex}%[1D2H3-2]%[ma:d4cd8a8d]
 Cho các dãy số ${a_n=n^2+n+1}$;${b_n=(n+2) \cdot 3^n}$; $\heva{&c_1=2 \\& c_{n+1}=\dfrac{6}{c_n}, \forall n \in \mathbb{N}^*}$ \break và ${d_n=(-4)^{2 n+1}}$.  Khi đó
  \choiceTF
  {\True ${\left(a_n\right)}$  không phải là cấp số nhân}
  {\True ${\left(b_n\right)}$ không phải  là cấp số nhân}
  {${\left(c_n\right)}$ là một cấp  số nhân}
  {\True ${\left(d_n\right)}$ là một cấp số  nhân }
  \loigiai{
  \begin{itemchoice}
   \itemch ${\dfrac{a_{n+1}}{a_n}=\dfrac{n^2+3 n+3}{n^2+n+1}, \forall n \in \mathbb{N}^*}$, không phải là hằng  số.\\
   Vậy  ${\left(a_n\right)}$ không phải là cấp số nhân.
   \itemch ${\dfrac{b_{n+1}}{b_n}=\dfrac{(n+3) \cdot 3^{n+1}}{(n+2) \cdot 3^n}=\dfrac{3(n+3)}{n+2}, \forall n \in \mathbb{N}^*}$, không phải là hằng số. Vậy ${\left(b_n\right)}$  không phải là cấp số nhân.
   \itemch Từ công thức truy  hồi của dãy số, suy ra ${c_1=2; c_2=3; c_3=2; c_4=3; \ldots}$.\\
   Vì ${\dfrac{c_3}{c_2} \neq  \dfrac{c_2}{c_1}}$ nên ${\left(c_n\right)}$ không phải là cấp số nhân.
   \itemch ${\dfrac{d_{n+1}}{d_n}=\dfrac{(-4)^{2(n+1)+1}}{(-4)^{2 n+1}}=16, \forall n  \in \mathbb{N}^*}$.\\
    Vậy ${\left(d_n\right)}$ là một cấp số nhân công bội ${q=16}$.
  \end{itemchoice}
 } 
\end{ex}

% nguồn: Ok_dungsai_lop11_Lan2
\begin{ex}%[1D2H3-2]%[ma:1a5c909b]
	Xác định tính đúng sai của các khẳng định sau 
	 \choiceTF
		{\True Dãy số $(u_n)$ với  $u_n = (-7)^n \cdot 5^{3n-1}$ là cấp số nhân với công bội $q = -875$ }
		{Dãy số $(u_n)$ với $\heva{& u_1  = 2 \\ & u_{n+1} = -5u_n}$ là cấp số nhân với công bội $q = -4$}
		{\True Dãy số $(u_n)$ với  $\heva{& u_1 = 2 \\ & u_{n+1} = u_n^2}$ không là cấp số nhân}
		{\True Dãy số $(u_n)$ với $-\dfrac{1}{8}$; $-\dfrac{1}{4}$ $-\dfrac{1}{2}$;  $1$ không là cấp số nhân}
	 \loigiai{
		\begin{itemchoice}
			\itemch Ta có 
			 \[
			\frac{u_{n+1}}{u_n} = \frac{(-7)^{n+1} \cdot 5^{3(n+1)-1}}{(-7)^n \cdot 5^{3n-1}} = (-7) \cdot  5^3 = -875.
			\] 
			Vậy  $(u_n)$ là cấp số nhân với công bội $q = -875$.
			 \itemch
			$\heva{& u_1 = 2 \\  & u_{n+1} = -5u_n}$. Khi đó:
			\[ 
			\frac{u_{n+1}}{u_n} = \frac{-5u_n}{u_n}  = -5.
			 \]
			Vậy $(u_n)$ là cấp số nhân với công bội $q = -5$ (không  phải $-4$).
			 \itemch
			Ta có $u_2 = u_1^2 = 2^2 = 4$; $u_3 = u_2^2 = 4^2  = 16$; $u_4 = u_3^2 = 16^2 = 256$.
			 Khi đó
			\[ 
			\frac{u_2}{u_1} = \frac{4}{2} = 2;  \quad \frac{u_4}{u_3} = \frac{256}{16} = 16.
			 \]
			Nhận thấy $\dfrac{u_2}{u_1} \ne \dfrac{u_4}{u_3}$  $(2 \ne 16)$.\\
			Vậy $(u_n)$ không là cấp số nhân. 
			\itemch 
			 \[
			\frac{u_2}{u_1} = \left(-\frac{1}{4}\right) \div \left(-\frac{1}{8}\right) = 2; \quad \frac{u_4}{u_3}  = 1 \div \left(-\frac{1}{2}\right) = -2
			\] 
			Nhận thấy: $\dfrac{u_2}{u_1} \ne \dfrac{u_4}{u_3}$, $(2  \ne -2)$.\\
			Vậy $(u_n)$  không là cấp số nhân.
		\end{itemchoice}
	 }
\end{ex}

% nguồn: Ok_dungsai_lop11_Lan2
\begin{ex}%[1D2H3-2]%[ma:12439e8a]
	Cho dãy số sau đây $u_n=\left(\sqrt{5}\right)^{2n-3}$; $v_n=\dfrac{2}{n}$; $w_n=\dfrac{3^{n+1}}{2^n}$ và dãy số hữu  hạn gồm các số hạng $16$; $4$; $1$; $\dfrac{1}{4}$; $\dfrac{1}{16}$; $\dfrac{1}{64}$. Xét tính đúng sai của các khẳng định sau
	 \choiceTF
		{\True $u_n$ là một cấp số nhân  công bội $q=5$}
		{\True $(v_n)$ không  phải là một cấp số nhân}
		{\True $(w_n)$ là một cấp số  nhân có số hạng đầu $w=\dfrac{9}{2}$}
		{ Dãy số hữu hạn đã cho theo thứ tự lập thành cấp số nhân có công  bội bằng $\dfrac{1}{8}$}
	\loigiai{ 
		\begin{itemchoice}
			 \itemch
			Ta có $u_{n+1}=\left(\sqrt{5}\right)^{2(n+1)-3}=\left(\sqrt{5}\right)^{2n-1}\Rightarrow \dfrac{u_{n+1}}{u_n}=\dfrac{\left(\sqrt{5}\right)^{2n-1}}{\left(\sqrt{5}\right)^{2n-3}}=\left(\sqrt{5}\right)^2=5$, $\forall x\in \mathbb{N}^*$.\\ 
			Do đó $(u_n)$ là một cấp  số nhân có số hạng đầu $u_1=\dfrac{1}{\sqrt{5}}$ với công bội $q=5$.
			\itemch 
			Ta có $v_{n+1} = \dfrac{2}{n+1} \Rightarrow \dfrac{v_{n+1}}{v_n} = \dfrac{n}{n+1}$ (tỉ số này còn phụ thuộc vào $n$). \\ 
			Do đó $(v_n)$ không phải là một cấp số nhân. 
			\itemch 
			Ta có 
			\[ 
			w_{n+1} = \dfrac{3^{n+2}}{2^{n+1}} = \dfrac{3}{2} \cdot \dfrac{3^{n+1}}{2^n} = \dfrac{3}{2} \cdot w_n \Rightarrow \dfrac{w_{n+1}}{w_n} = \dfrac{3}{2} ,\quad \forall n \in  \mathbb{N}^*.
			 \]
			Do đó $(w_n)$ là một cấp số nhân có số hạng đầu $w_1 = \dfrac{3^2}{2^1} = \dfrac{9}{2}$ với công bội $q =  \dfrac{3}{2}$.
			 \itemch
			Đặt $u_1 = 16$ , $u_2 = 4$ , $u_3 = 1$; $u_4 = \dfrac{1}{4}$ , $u_5 =  \dfrac{1}{16}$ , $u_6 = \dfrac{1}{64}$. \\
			Ta có 
			\[ 
			u_2 = \dfrac{u_1}{4} ,\quad u_3 = \dfrac{u_2}{4} ,\quad u_4 = \dfrac{u_3}{4} ,\quad  u_5 = \dfrac{u_4}{4} ,\quad u_6 = \dfrac{u_5}{4}.
			\] 
			Vì vậy dãy số  hữu hạn đã cho theo thứ tự lập thành cấp số nhân có công bội bằng $\dfrac{1}{4}$.
		\end{itemchoice}
	 }
\end{ex}

% nguồn: Ok_dungsai_lop11_Lan2
\begin{ex}%[1D2H3-2]%[ma:19beb3d2]
	Cho các dãy số $a_n=n^2+n+1$; $b_n=(n+2)\cdot 3^n$; $\left\{\begin{array}{l}{c_1=2} \\ {c_{n+1}=\dfrac{6}{c_n}, \forall n\in \mathbb{N}^{*}}
\end{array}\right.$; $d_n=(-4)^{2n+1}$. Khi đó
 \choiceTF
	{\True $\left(a_n \right)$ không phải là  cấp số nhân}
	{\True $\left(b_n \right)$ không phải  là cấp số nhân}
	{$\left(c_n \right)$ là một cấp số  nhân}
	{\True $\left(d_n \right)$  là một cấp số nhân}
 \loigiai{
	\begin{itemchoice}
		\itemch \textbf{Đúng.}  \\
		$\dfrac{a_{n+1}}{a_n}=\dfrac{n^2+3n+3}{n^2+n+1}, \forall n\in \mathbb{N}^{*}$, không phải là hằng số. Vậy $\left(a_n \right)$ không phải là cấp số nhân. 
		 \itemch \textbf{Đúng.} \\
		$\dfrac{b_{n+1}}{b_n}=\dfrac{(n+3)\cdot 3^{n+1}}{(n+2)\cdot 3^n}=\dfrac{3(n+3)}{n+2}, \forall n\in  \mathbb{N}^{*}$, không phải là hằng số. Vậy $\left(b_n \right)$ không phải là cấp số nhân
		\itemch  \textbf{Sai.} \\
		Từ công thức truy hồi của  dãy số, suy ra $c_1=2;c_2=3;c_3=2;c_4=3;\ldots$
		Vì $\dfrac{c_3}{c_2} \ne \dfrac{c_2}{c_1}$ nên  $\left(c_n \right)$ không phải là cấp số nhân
		\itemch \textbf{Đúng.}\\ 
		$\dfrac{d_{n+1}}{d_n}=\dfrac{(-4)^{2(n+1)+1}}{(-4)^{2n+1}}=16,\forall n\in \mathbb{N}^{*}$. Vậy $\left(d_n  \right)$ là một cấp số nhân công bội $q=16$.
	\end{itemchoice}
 }
\end{ex}

% nguồn: Ok_dungsai_lop11_Lan1
\begin{ex}%[1D2C3-2]%[ma:54b5f082]
 Cho cấp số nhân $\left(u_n\right)$ có các số hạng thỏa mãn  $\heva{&u_2=8\\&3u_1+2u_2-u_3=20.}$
 \choiceTF 
  {\True Tồn tại hai cấp số nhân $\left(u_n\right)$ thoả mãn} 
  {Cấp số  nhân $\left(u_n\right)$ có công bội là $-3$ hoặc $\dfrac{3}{2}$}
  {\True Số hạng $u_5$ của cấp số  nhân trên với công bội $q<0$ là $-64$}
  {\True Tổng $20$ đầu của cấp  số nhân là $1\,398\,100$}
 \loigiai{ 
  \begin{eqnarray*}
    & &\heva{&u_2=8\\&3u_1+2u_2-u_3=20}
   \Leftrightarrow \heva{&u_2=qu_1=8\\&3u_1+2\cdot 8-8q=20} \Leftrightarrow \heva{&u_1=\dfrac{8}{q}\\&3\cdot  \dfrac{8}{q}+16-8q=20} \\
   &\Leftrightarrow&  \heva{&u_1=\dfrac{8}{q}\\&8q^2+4q-24=0}
   \Leftrightarrow \heva{&u_1=-4\\&\hoac{&q=-2\\&q=\dfrac{3}{2}.}} 
  \end{eqnarray*}
  \begin{itemchoice}
   \itemch  {\bf Đúng}. \\
   Với $q=\dfrac{3}{2}$ hoặc $q=-2$,  ta được hai cấp số nhân $\left(u_n\right)$ thoả mãn.
    \itemch {\bf Sai}. \\
   Cấp số nhân $\left(u_n\right)$ có công bội là  $-2$ hoặc $\dfrac{3}{2}$.
   \itemch {\bf  Đúng}. \\
   Ta có  $u_5=u_1\cdot q^4=-4\cdot (-2)^4=-64$.
   \itemch  {\bf Đúng}. \\
   Ta  có $S_{20}=\dfrac{u_1\left(1-q^{20}\right)}{1-q}$.\\
   Với $q=-2$ thì $S_{20}=1\,398\,100$. 
  \end{itemchoice}
 } 
\end{ex}

% nguồn: Ok_dungsai_lop11_Lan2
\begin{ex}%[1D2C3-2]%[ma:9c963a8e]
	Cho  các dãy số $a_n = n^2 + n + 1$; $b_n = (n + 2) \cdot 3^n$;
	$\heva{&c_1 = 2 \\&c_{n+1} = \dfrac{6}{c_n},\ \forall n \in \mathbb{N}^*}$; $d_n  = (-4)^{2n+1}$.
	Xét tính đúng sai của  các khẳng định sau
	 \choiceTF
		{\True  $\left(a_n\right)$ không phải là cấp số nhân}
		 {\True $\left(b_n\right)$ không phải là cấp số nhân}
		 {$\left(c_n\right)$ là một cấp số nhân}
		 {\True $\left(d_n\right)$ là một cấp số nhân}
	 \loigiai{
		\begin{itemchoice}
			 \itemch
			$\dfrac{a_{n+1}}{a_n} = \dfrac{n^2 +  3n + 3}{n^2 + n + 1}$, $\forall n \in \mathbb{N}^*$. Không phải là hằng số nên $(a_n)$ không phải là cấp số nhân.
			\itemch 
			$\dfrac{b_{n+1}}{b_n} = \dfrac{(n+3)\cdot 3^{n+1}}{(n+2)\cdot 3^n} = \dfrac{3(n+3)}{n+2}$, $\forall n \in \mathbb{N}^*$. Không phải là hằng số nên $(b_n)$ không phải là cấp số nhân. 
			\itemch 
			Xét công thức truy hồi của dãy $(c_n)$, ta có 
			$\heva{&c_1 = 2 \\&c_2 = \dfrac{6}{2} =  3 \\&c_3 = \dfrac{6}{3} = 2 \\&c_4 = \dfrac{6}{2} = 3 \\&\ldots}$
			 \(\Rightarrow (c_n)\) là dãy xen kẽ $2,\ 3,\ 2,\ 3,\ \ldots$ \\
			Vì $\dfrac{c_3}{c_2}\neq \dfrac{c_2}{c_1}$ nên $(c_n)$ không  phải là cấp số nhân.
			 \itemch
			$\dfrac{d_{n+1}}{d_n} = \dfrac{(-4)^{2(n+1)+1}}{(-4)^{2n+1}} = (-4)^2  = 16$, $\forall n \in \mathbb{N}^*$. \\
			Vậy $(d_n)$ là một cấp số nhân với  công bội $q = 16$.
		\end{itemchoice}
	} 
\end{ex}

% nguồn: Ok_dungsai_lop11_Lan1
\begin{ex}%[1D2B3-2]%[ma:b51ef36d]
 Hãy cho biết tính đúng sai của mỗi mệnh đề dưới  đây?
  \choiceTF
  {Cho cấp số nhân $\left(u_n \right)$, biết $u_1=200$ và $u_2=-800$. Công bội của cấp số nhân đã cho là  $q=4$}
  {Cấp số nhân $(u_n)$ có $u_5=6$, $u_6=2$. Công bội của cấp số nhân đó bằng  $3$}
  {Cho cấp số nhân $\left(u_n\right)$, biết $u_1=-2$. Tổng của hai số hạng đầu $S_2=6$. Cấp số nhân đó có công  bội bằng $4$}
  {\True Cho cấp số nhân $(u_n)$ thỏa mãn $\heva{&u_{10}=8u_7 \\&u_1+u_4=144}$. Công bội $q$ của cấp  số nhân $(u_n)$ bằng $3$}
  \loigiai{
  \begin{itemchoice}
    \itemch Ta có $q=\dfrac{u_2}{u_1}=-4$.
   \itemch  Ta có $q=\dfrac{u_6}{u_5}=\dfrac{1}{3}$.
   \itemch Ta có  $S_2 = u_1 + u_2 \Rightarrow u_2 = S_2 - u_1 = 8$. \\
   Vậy công bội  $q = \dfrac{u_2}{u_1} = -4$.
    \itemch Ta có $u_{10}=8u_{7} \Leftrightarrow u_7\cdot q^3 =8u_7 \Leftrightarrow \hoac{&u_7=0\\&q=3.}$\\
   Vậy  $q=3$.
  \end{itemchoice}
 } 
\end{ex}

% nguồn: Ok_dungsai_lop11_Lan1
\begin{ex}%[1D2B3-2]%[ma:f7b67a2f]
 Hãy cho biết tính đúng sai của mỗi mệnh  đề dưới đây?
 \choiceTF 
  {\True Dãy số $1$, $-2$,  $4$, $-8$, $16$ là một cấp số nhân}
  {\True Dãy số $\heva{&u_1=-1\\&u_{n+1}=-3u_n,n\geq 1}$ là một cấp  số nhân}
  {Dãy số $u_n=(-1)^n\cdot2$ không  là một cấp số nhân}
  {Dãy số $u_n=2^n+2$  là một cấp số nhân}
  \loigiai{
  \begin{itemchoice}
   \itemch Ta có $1,-2,4,-8,16$ là cấp số nhân có số hạng đầu $u_1=1$,  công bội $q=-2$.
   \itemch Dãy số $\heva{&u_1=-1\\&u_{n+1}=-3u_n,n\geq 1}$ là một cấp số nhân có công bội $q=\dfrac{u_{n+1}}{u_n}=-3$. 
   \itemch Đặt $(u_n)_{n\ge 1}$ là dãy số $-2;2;-2;2; \ldots$\\ 
   Ta có $  u_{n+1}=-u_n$, $ \forall n\ge 1 $ nên $(u_n)$ là cấp số nhân.
   \itemch Xét dãy số $u_n=2^n+2$, có $u_1=4, u_2=6, u_3=10$.\\ 
   Suy ra $\dfrac{u_3}{u_2}=\dfrac{5}{3}\neq \dfrac{u_2}{u_1}=\dfrac{3}{2}$.\\ 
   Vậy  $u_n=2^n+2$ không phải cấp số nhân.
  \end{itemchoice}
  }
\end{ex}


%%%%%%%%%%%%%%%% TRẢ LỜI NGẮN %%%%%%%%%%%%%%%%

% nguồn: Ok_TraLoiNgan_Lop11_Lan2
\begin{ex}%[1D2V3-2]%[ma:d74f3058]
	Một cấp số nhân có $6$ số hạng, số hạng đầu bằng $2$ và số hạng thứ $6$ bằng $486$. Tìm công bội $q$ của cấp số  nhân đã cho.
	 \shortans{3}
	 \loigiai{
		Theo giả thiết  ta có
		\[\heva{&u_1=2\\&u_6=486}  \Rightarrow 486=u_6=u_1q^5=2q^5 \Leftrightarrow q^5=243\Leftrightarrow q=3.\]
	} 
\end{ex}


%%%%%%%%%%%%%%%% TỰ LUẬN %%%%%%%%%%%%%%%%

% nguồn: Ok_TuanLuan_Toan11_Lan2
\begin{ex}%[1D2H3-2]%[ma:fc8a47da]
	Xét trong các dãy số  sau, dãy số nào là cấp số nhân, (nếu có) và tìm công bội của cấp số nhân đó.
	\begin{listEX}[2]
		\item $u_n  = (-3)^{2n+1}$.
		\item $u_n = (-1)^n  \cdot 5^{3n+2}$.
		\item $\heva{&u_1 = 2  \\ &u_{n+1} = u_n^2.}$
		\item  $\heva{&u_1 = 3 \\ &u_{n+1} = \dfrac{9}{u_n}.}$
	\end{listEX}
	\loigiai{ 
		\begin{listEX}
			 \item Ta có
			\[ 
			\dfrac{u_{n+1}}{u_n}  = \dfrac{(-3)^{2n+3}}{(-3)^{2n+1}} = (-3)^2 = 9 \quad \text{(không đổi)}
			\] 
			nên $(u_n)$ là cấp số  nhân với công bội $q = 9$.
			\item Ta có 
			 \[
			\dfrac{u_{n+1}}{u_n} = \dfrac{(-1)^{n+1} \cdot 5^{3(n+1)+2}}{(-1)^n \cdot 5^{3n+2}} = (-1)^1 \cdot  5^3 = -125 \quad \text{(không đổi)}\]
			 nên $(u_n)$ là cấp số nhân với công bội $q = -125$.
			\item Ta có $u_2 = u_1^2 = 4$, $u_3 =  u_2^2 = 16$, $u_4 = u_3^2 = 256$.\\
			Suy ra $\dfrac{u_2}{u_1} = \dfrac{4}{2} = 2$, và $\dfrac{u_4}{u_3} = \dfrac{256}{16}  = 16$\\
			$\Rightarrow \dfrac{u_2}{u_1} \ne \dfrac{u_4}{u_3}$, do đó $(u_n)$ không  là cấp số nhân.
			\item Ta có $u_{n+1} = \dfrac{9}{u_n}$ và $u_{n}  = \dfrac{9}{u_{n-1}}$ với mọi $n \ge 2$. Khi đó
			 \[
			u_{n+1}=\dfrac{9}{\dfrac{9}{u_{n-1}}} = u_{n-1},~\forall n  \ge 2.
			\] 
			Do đó,  ta có
			\[u_1 = u_3 = u_5  = \ldots = u_{2n+1} \quad (1)\]
			và 
			\[ u_2 = u_4 = u_6 = \ldots  = u_{2n} = \ldots \quad (2)\]
			Theo đề bài $u_1 = 3 \Rightarrow u_2 = \dfrac{9}{u_1}  = \dfrac{9}{3} = 3$. \quad (3)
			Từ (1),  (2) và (3) suy ra $u_1 = u_2 = u_3 = u_4 =u_5= \ldots = u_{2n} =u_{2n+1} \ldots$.\\
			Kết luận $(u_n)$  là cấp số nhân với công bội $q = 1$.
		\end{listEX}
	 }
\end{ex}

% nguồn: Ok_TuanLuan_Toan11_Lan2
\begin{ex}%[1D2H3-2]%[ma:5137316c]
	Giá  trị của $a$ để $\dfrac{-1}{5}$; $a$; $\dfrac{-1}{125}$ theo thứ tự lập thành cấp số nhân?
	\loigiai{ 
		Ta có $a^2=\left(-\dfrac{1}{5}\right)  \cdot\left(-\dfrac{1}{125}\right)=\dfrac{1}{625} \Leftrightarrow a=\pm \dfrac{1}{25}$.
	 }
\end{ex}

% nguồn: Ok_TuanLuan_Toan11_Lan2
\begin{ex}%[1D2H3-2]%[ma:609d6b9a]
	Chứng minh rằng dãy số  $\left(v_n\right): v_n=(-1)^n \cdot 3^{2 n}$ là một cấp số nhân.
	 \loigiai{
		 Ta có $\dfrac{v_{n+1}}{v_n}=\dfrac{(-1)^{n+1}\cdot 3^{2(n+1)}}{(-1)^n\cdot 3^{2 n}}=-9, \forall n \in \mathbb{N}^*$.\\
		Vậy $\left(v_n\right):  v_n=(-1)^n \cdot 3^{2 n}$ là một cấp số nhân.
	} 
\end{ex}

% nguồn: Ok_TuanLuan_Toan11_Lan2
\begin{ex}%[1D2V3-2]%[ma:4ea527f3]
	Cho dãy số $\left(u_n\right)$ được xác định bởi $\heva{&u_1=2\\&u_{n+1}=4u_n+9}$, $\forall n\geq 1$. \\ 
	Chứng minh rằng dãy số $\left(v_n\right)$ xác định bởi $v_n=u_n+3$, $\forall n\geq 1$ là một cấp số nhân.  Hãy xác định số hạng đầu và công bội của cấp số nhân đó.
	\loigiai{ 
		Ta có $v_n = u_n + 3$ \quad (1) $\Rightarrow v_{n+1} =  u_{n+1} + 3$ \quad (2).\\
		 Theo đề $u_{n+1} = 4u_n + 9 \Rightarrow u_{n+1} + 3 = 4(u_n + 3)$ \quad (3).\\
		Thay (1) và (2) vào (3) được  $v_{n+1} = 4v_n$, $\forall n \geq 1 \Rightarrow \dfrac{v_{n+1}}{v_n} = 4$ (không đổi).\\
		Kết luận $\left(v_n\right)$ là cấp số nhân với công bội $q = 4$ và số hạng đầu $v_1 = u_1  + 3 = 5$.
	 }
\end{ex}
